% s14 The coupling range
\section{The coupling range}

\providecommand{\sxivNE}{\ensuremath{(\mathrm{NE})}}
\providecommand{\sxivNEl}{\ensuremath{(\mathrm{NE}^{<})}}
\providecommand{\sxivCT}{\ensuremath{(\mathrm{CT}_{\infty})}}

This section collects what the sources behind the inputs of \cref{thm:s13:main} say about the couplings and shows that the only coupling condition anywhere in this paper is the one-sided sign clause \Tthree\ of the Wick leg, which \cref{prop:s14:T3} proves on explicit regions.

\medskip\noindent\textbf{King's range, units, two letters $J$.} Let $e>0$ and $\lambda,m,\lmu>0$. Every constant of the cited constructions, including the number $J_{\rm K}$ of renormalisation steps after which the running couplings are small and the stopping scale $5^{-J_{\rm K}}$ (a member of the family $\{5^{-k}\}$), may depend on $(e,\lambda,m,\lmu)$, on the large-block integer $M$ and on the cube or torus; no numerical restriction is placed on the couplings themselves. That a stopping scale exists at every positive point is part of the cited constructions (\cite[Thm.~3.1]{K1}, \cite{Ba1,Ba2,Ba3,Ba4}, and the existence and rotation parts of this paper), under the mass cap of \cite[Prop.~3.1]{Ba2}, a one-sided cap on the stopping scale (\cref{rem:s14:BALm}). This is what ``King's range'' means here: all positive $(e,\lambda,m,\lmu)$, with a stopping scale chosen after them. The cube index $J$ of $\CJ$ is ours; King's stopping index, the ``fixed integer $J$'' of \cite[Thm.~3.1]{K1}, is written $J_{\rm K}$. Couplings are measured in the units in which the coarsest spacing is $1$; in $d=3$ all four of $e^2,\lambda,m,\lmu$ have the dimension of a mass, rescaling to the unit lattice at spacing $\eps$ replacing $m^2,\lambda,e$ by $m^2\eps^2,\lambda\eps,e\eps^{1/2}$ \cite[Sect.~0]{K1} and $\lmu^2$ like $m^2$.

\medskip\noindent\textbf{The relative bare coefficient.} Let $\alpha^\infty_k(M)=M^2+\dm_k(M)$, $s^\infty_k(M_1)=\frac12(\alpha^\infty_k(m)-\alpha^\infty_k(M_1))$ and $\sigma_1(M_1)=\liminf_ks^\infty_k(M_1)$ be as in \cref{hyp:s03:T2}; the theorem is stated through $\sigma_1$ alone. Describing \emph{where} \Tthree\ fails uses one further statement about the continuum limit of this coefficient, which is of the same kind as \cref{hyp:s03:T2} (\cref{rem:s03:T2}) and which we name rather than use tacitly:

\smallskip\noindent\sxivCT\ \emph{For every $M>0$ the limit
\[
\alpha_\infty(M)-\alpha_\infty(m)\defeq\lim_{k\to\infty}\big[\alpha^\infty_k(M)-\alpha^\infty_k(m)\big]
\]
exists and is a polynomial of degree $\le2$ in $(e^2,\lambda)$ with finite coefficients, differentiable in $M^2$ at $M=m$.}

\smallskip\noindent The divergent parts of the counterterm ($1/\eps$ at one loop, $\ln(1/\eps)$ at two loops) do not depend on the scalar mass and cancel in the difference. \sxivCT\ is proved in \cref{thm:a06:CD1}(b); it is not a hypothesis of \cref{thm:s13:main} and enters only the description of the failure set. Under \sxivCT\ put
\[
s_\infty(M_1)\defeq\tfrac12\big(\alpha_\infty(m)-\alpha_\infty(M_1)\big),\qquad s_*\defeq\sup_{M>0}s_\infty(M)\in[0,\infty],
\]
\[
E_W\defeq\{s_*=0\}=\{m\text{ is a global minimiser of }M\mapsto\alpha_\infty(M)\text{ on }(0,\infty)\},\qquad
E^<_W\defeq\{m\text{ minimises }\alpha_\infty\text{ on }(0,m]\}\supseteq E_W,
\]
subsets of the coupling space ($s_*\ge0$ because $M=m$ gives $0$; the value $+\infty$ is harmless below). The two sign clauses are
\begin{itemize}[nosep,leftmargin=*]
\item[\sxivNE] $m$ is not a global minimiser of $\alpha_\infty$ on $(0,\infty)$, i.e.\ $s_*>0$, i.e.\ $(e,\lambda,m,\lmu)\notin E_W$;
\item[\sxivNEl] some $M_1\in(0,m)$ has $\alpha_\infty(M_1)<\alpha_\infty(m)$, i.e.\ $(e,\lambda,m,\lmu)\notin E^<_W$ (the one-sided form; it implies \sxivNE).
\end{itemize}
Throughout, $'$ denotes $d/dM^2$, $N$ is the number of real components of the scalar field, and ``one-loop order'' means the orders $e^2$ and $\lambda$.

\begin{lemma}[the coupling range]\label{lem:s14:CW}
Let $e>0$ and $\lambda,m,\lmu>0$.
\begin{enumerate}[label=\textup{(\alph*)},leftmargin=*]
\item \emph{The curvature sector.} The existence theorem of the curvature sector, \cref{hyp:s03:EC} (the existence clause of \cite[Thm.~2.1]{K2} at zero scalar source), carries no condition on $e,\lambda,m,\lmu$. The correlation inequalities it uses, \cref{lem:s02:BFS} (Theorem~6.2(1),(2) of \cite{BFS1}), have the standing hypotheses of a finite single-site measure, positive hopping coefficients and a positive photon covariance, which for King's action are $\lambda>0$ and $\lmu>0$; the margin $R(\lmu)=\pi\sqrt3/(2\lmu)$ of \cref{lem:s07:U} needs $\lmu>0$ only; and no lemma proved in this paper for the curvature sector carries a condition on the couplings beyond those of its inputs. Hence the curvature sector of \cref{thm:s13:main} holds on all of King's range, with no exceptional set. The one proof-internal proviso in the Ba{\l}aban antecedents of \cref{hyp:s03:EC}, the mass clause of \cite[Prop.~3.1]{Ba2}, is a cap on the dimensionless scalar mass at the stopping scales and no condition on the bare couplings (\cref{rem:s14:BALm}). Of the bounds of the paper only the absolute bounds of \cref{thm:s05:CK} are uniform in $\lmu$.
\item \emph{The Wick sector, Form~(B).} The inputs of Form~(B) beyond \cref{hyp:s03:W1,hyp:s03:T1,hyp:s03:T2} and \Tthree\ --- \cref{prop:s10:RP,lem:s10:CB,lem:s02:BFS,lem:s10:F,lem:s10:V} and the tube statements \cref{thm:s11:W,lem:s11:TWtube,thm:s08:OS2WT}, whose conditions on the Wick amplitudes are not conditions on the couplings --- print no condition on the bare couplings beyond positivity ($\lambda>0$, $\lmu>0$); \cref{hyp:s03:T1} is used at the two masses $m$ and $M_1$. The only coupling condition of Form~(B) is therefore \Tthree, which \cref{prop:s14:T3} proves on explicit regions. Under \sxivCT:
\begin{enumerate}[label=\textup{(\roman*)},nosep,leftmargin=*]
\item $\sigma_1(M_1)=\lim_ks^\infty_k(M_1)=s_\infty(M_1)$ for every $M_1>0$. Hence \Tthree\ holds if and only if \sxivNEl: it fails exactly on $E^<_W$. If \cref{hyp:s03:T1,hyp:s03:T2} are available also at auxiliary masses $M_1>m$, then by \cref{prop:s10:Upm} the two-sided clause \sxivNE\ replaces \Tthree, with failure set exactly $E_W$; the branch $M_1>m$ is needed exactly on $E^<_W\setminus E_W$, for instance when $\alpha'_\infty(m)<0$ and $m$ minimises $\alpha_\infty$ on $(0,m]$.
\item To one-loop order, for every $M_1>0$,
\[
s_\infty(M_1)=\tfrac12(m^2-M_1^2)+\frac{\lambda(N+2)(m-M_1)}{2\pi}-e^2\,\frac{(m-M_1)\,(m\lmu+mM_1+M_1\lmu)}{8\pi(\lmu+M_1)(m+\lmu)},
\]
so that
\begin{equation}\label{eq:s14:alphaprime}
\alpha'_\infty(m)=1+\frac{\lambda(N+2)}{2\pi m}-e^2\,\frac{m+2\lmu}{8\pi(m+\lmu)^2},\qquad \frac{m+2\lmu}{8\pi(m+\lmu)^2}\le\frac1{4\pi(m+\lmu)},
\end{equation}
and to one-loop order \Tthree\ holds if and only if
\begin{equation}\label{eq:s14:ec2}
e^2<e_c^2(\lambda,m,\lmu)\defeq\max\Big\{\frac{8\pi(m+\lmu)^2}{m+2\lmu}\Big(1+\frac{\lambda(N+2)}{2\pi m}\Big),\ 4\pi(m+\lmu)\Big(1+\frac{\lambda(N+2)}{\pi m}\Big)\Big\},
\end{equation}
the first entry being the larger if and only if $\lambda(N+2)\lmu\le\pi m^2$. So the one-loop failure set of \Tthree\ is the half-line $e^2\ge e_c^2$. The first entry is the zero $e_*^2\ge4\pi(m+\lmu)$ of \eqref{eq:s14:alphaprime}; at $m=\lmu=1$, $\lambda=0$ it equals $32\pi/3\approx33.5$ (for orientation only). Nothing finer than one-loop order is claimed about $E^<_W$.
\item For the two-sided set, $E_W\subset\{\alpha'_\infty(m)=0\}$; for each $(\lambda,m,\lmu)$ at which the polynomial $e^2\mapsto\alpha'_\infty(m)$, of degree $\le2$, is not identically zero, this is at most two values of $e^2$. At $\lmu=m$ the window $0<\lambda,e^2<0.1m$ of the torus paper \cite{SymPhase} lies outside $E^<_W$, by (d).
\end{enumerate}
\item \emph{The existence and rotation parts.} The existence theorems \cref{hyp:s03:EC,hyp:s03:W1} and the rotation theorem \cref{hyp:s03:RC} are proved for $e>0$, $\lambda,m,\lmu>0$ in King's range, and none carries a condition on the bare couplings beyond positivity: every smallness condition of \cite{Ba1,Ba2,Ba3,Ba4} that they use (listed in the proof) is a one-sided cap on a stopping scale $\beta\le1$, imposed last in the schedule of constants, and an admissible $\beta$ exists at every positive point.\item \emph{Comparison with the torus paper.} The window $0<\lambda<0.1m$, $0<e^2<0.1m$, $\lmu\in(0,m]$ of \cite{SymPhase} is a hypothesis of its torus theorem only and is not inherited by King's state. Nor is $\lmu\le m$: the one King-side input printed under it, the torus paper's reflection-positivity proposition behind \cref{prop:s10:RP} (stated for $\lmu\in(0,\mus]$, $\mus=m$), is lifted to all $\lmu>0$ by a proof that uses only site-independent coefficients, $\lambda>0$ and even sides. At $\lmu=m$ the torus paper's certificate gives $s_\infty(0.7m)\ge0.1595\,m^2$ on the whole window (identifying the two torus counterterm schemes at $\lmu=m$, which compute the same graphs), so \Tthree\ holds there with $M_1=0.7m$ and an explicit margin; \cref{prop:s14:T3} sharpens this and gives explicit regions at every $\lmu\in(0,m]$ (\cref{rem:s14:torus-window}).
\item \emph{The quantifier.} For all $e>0$, $\lambda,m,\lmu>0$ and each datum there are an integer $J_{\rm K}$ (\cite[Thm.~3.1]{K1}) and a stopping scale $\beta\le1$ such that the running couplings $(e\beta^{1/2},\lambda\beta,m\beta,\lmu\beta)$ satisfy finitely many one-sided caps on $\beta$ and on three of the four of them ($\lmu\beta$ is uncapped in the sources); no condition is placed on the bare couplings, and the theorems consumed from \cite{K1,K2} carry no coupling hypothesis. A sentence ``$e,\lambda$ sufficiently small, $m,\lmu>0$ fixed'', read as a hypothesis on the bare couplings, is strictly stronger and is printed in none of the sources (\cref{rem:s14:quantifier}).
\end{enumerate}
\end{lemma}

\begin{proof}
\emph{(a).} Theorem~2.1 of \cite{K2} asserts, for $g,h$ supported in $C_J$ and all $z,w\in\C$, the existence of $\lim_K\langle\Theta(zg,wh)\rangle^{\eps_K}_{N,D,C_J}$; no hypothesis on the couplings precedes it. Its proof rests on \cite[Thm.~3.1]{K1} (the convergence of the effective action, referred there to \cite{Ba1,Ba2,Ba3,Ba4}), stated ``for some fixed integer $J$'', our $J_{\rm K}$; neither it nor \cite[Thms.~2.1, 3.4]{K1} carries a coupling hypothesis, the smallness in \cite{K1} being of the fields and of the stopping spacing. So the one proviso of \cite[Prop.~3.1]{Ba2} is upstream of \cref{hyp:s03:EC}. For King's action the single-site measure $e^{-\eps^3[\frac12\alpha r^2+\lambda r^4]}r\,dr$ is finite iff $\lambda>0$ (or $\lambda=0$, $\alpha>0$), the hopping weights are positive, and the Feynman-form photon covariance is positive iff $\lmu>0$; the ``$g>0$'' of Theorem~6.2 of \cite{BFS1} is the sign of the scalar source, and \cref{lem:s02:BFS} covers $g=0$. The margin of \cref{lem:s07:U} concerns the photon Gaussians of the two boxes alone, so $\lmu$ is its only parameter (no fixed margin works at $\lmu=0$). The lemmas of this paper see only the form of the action.

\emph{The torus inputs.} \Cref{hyp:s03:T1} is Ba{\l}aban's stability theorem \cite[(1.14)]{Ba1}, \cite[(1.3)]{Ba2}, with constants independent of $\eps$ and of the torus, under $\lambda>0$, $\lmu^2>0$, $m^2>0$. Every printed smallness in \cite{Ba1,Ba2,Ba3,Ba4} is of one of three kinds. \emph{Running charge:} the bounds hold ``for $e(L^k\eps)$ sufficiently small'', $e(L^k\eps)=e\,(L^k\eps)^{(4-d)/2}$ \cite[(2.23)]{Ba1}; \cite{Ba4} proves them for a unit-lattice model whose charge is this running charge (its theorem is stated there as Proposition~2.1 of \cite{Ba1}), so ``$e$ sufficiently small'' there is a cap on $e\beta^{1/2}$ at the stopping scale $\beta$. \emph{Scale:} the stopping scale $\eps_0$ with $L^K\eps\le\eps_0<L^{K+1}\eps$ is ``some fixed sufficiently small number'' \cite{Ba3}, and ``$L^{k+1}\eps\le1$'', ``$L^k\eps$ sufficiently small'' \cite{Ba1,Ba2} are caps on it; Ba{\l}aban's units have $\lambda=1$ (in $d=3$: spacing $\lambda\eps$, mass squared $m^2/\lambda^2$, charge $e\lambda^{-1/2}$, the running charge invariant), so his ``$\eps$ small'' is ``$\lambda\eps$ small'' in ours, and his constants may depend on $e\lambda^{-1/2}$, $m/\lambda$, $\lmu/\lambda$: a dependence, not a condition. \emph{Mass:} Proposition~3.1 of \cite{Ba2} carries the sentence ``we assume $m^2\le O(1)$ also''; this is the mass cap of that proposition (\cref{rem:s14:BALm}). \Cref{hyp:s03:T2} carries no coupling condition: the bound on the relative bare coefficient holds for every mass in a compact interval containing $m$, on either side of $m$, with a constant depending on the couplings (\cref{rem:s03:T2}; the power-counting theorem of \cite{Rei88} that enters it has none). \Cref{lem:s10:F} needs a finite single-site integral ($\lambda>0$) and a finite photon Gaussian on the torus ($\lmu>0$; at $\lmu=0$ the constant mode lies in the kernel); \cref{lem:s10:V} needs massive propagators; \cref{prop:s10:RP} needs finite single-site moments and a positive Proca covariance and does not see the counterterm. The tube statements restrict Wick amplitudes only. So the only coupling condition of Form~(B) is the sign clause of \cref{prop:s10:U}, i.e.\ \Tthree.

\emph{(b)(i).} By \cref{lem:s10:V}, $s^\infty_k(M_1)$ is the fixed-$k$ volume limit of the torus coefficient, the relative bare coefficient on $\Lk$ in the free convention; by \sxivCT\ its limit $s_\infty(M_1)$ exists, hence $\sigma_1(M_1)=\liminf_k=\lim_k=s_\infty(M_1)$. Thus some $M_1\in(0,m)$ has $\sigma_1(M_1)>0$ iff \sxivNEl, and some $M_1>0$ has $\sigma_1(M_1)>0$ iff \sxivNE. \Cref{prop:s10:Upm} gives \Itwo\ from \cref{hyp:s03:T1,hyp:s03:T2} at an auxiliary mass on either side of $m$, so with those inputs \sxivNE\ suffices; on $E^<_W\setminus E_W$ every $M_1<m$ fails and some $M_1>m$ works.

\emph{(b)(ii).} The closed form is the one-loop evaluation of King's counterterm \cite[(3.32)--(3.33)]{K1} (tadpole and zero-momentum photon-exchange graphs; the $1/\eps$ divergence cancels in the difference). It equals $\frac12\int_{M_1^2}^{m^2}\alpha'_\infty(\mu)\,d\mu^2$ with $\alpha'_\infty(\mu)=1+\frac{\lambda(N+2)}{2\pi\mu}-e^2\frac{\mu+2\lmu}{8\pi(\mu+\lmu)^2}$, an orientation-free definite integral, valid on either side of $m$; at $\mu=m$ this is \eqref{eq:s14:alphaprime}, whose inequality is $m\ge0$. For $M_1\in(0,m)$ write $s_\infty(M_1)=\frac{m-M_1}2\beta(M_1)$ with
\[
\beta(M_1)=(m+M_1)+c-e^2\,\frac{m\lmu+(m+\lmu)M_1}{4\pi(m+\lmu)(\lmu+M_1)},\qquad c\defeq\frac{\lambda(N+2)}{\pi};
\]
$s_\infty(M_1)>0$ iff $\beta(M_1)>0$ iff $e^2<e_c^2(M_1)\defeq\big[(m+M_1)+c\big]\frac{4\pi(m+\lmu)(\lmu+M_1)}{m\lmu+(m+\lmu)M_1}$. With $\mathcal N(M_1)\defeq(m+M_1+c)(\lmu+M_1)$ and $\mathcal D(M_1)\defeq m\lmu+(m+\lmu)M_1$, the derivative of $e_c^2(M_1)$ has the sign of $\mathcal N'\mathcal D-\mathcal N\mathcal D'=2m\lmu M_1+(m+\lmu)M_1^2-c\lmu^2$, positive on $(0,m)$ when $c=0$ and, when $c>0$, negative then positive with one zero. So $e_c^2(\cdot)$ is increasing, or decreasing then increasing, on $(0,m)$, and its supremum over the open interval is the larger of the two end values, not attained: $M_1\to m^-$ gives $\frac{8\pi(m+\lmu)^2}{m+2\lmu}(1+\frac c{2m})$, the zero $e_*^2$ of \eqref{eq:s14:alphaprime}, and $M_1\to0^+$ gives $4\pi(m+\lmu)(1+\frac cm)$. Hence to one-loop order \Tthree\ holds iff $e^2<e_c^2(\lambda,m,\lmu)$ as in \eqref{eq:s14:ec2}, and fails for every $e^2\ge e_c^2$ (then $\beta\le0$ for every $M_1$); the first entry is the larger iff $(m+\lmu)(2m+c)\ge(m+c)(m+2\lmu)$, i.e.\ $c\lmu\le m^2$.

\emph{(b)(iii).} Under \sxivCT, $M^2\mapsto\alpha_\infty(M)$ is differentiable at $m$, and a global minimum at the interior point $m$ is a critical point: $E_W\subset\{\alpha'_\infty(m)=0\}$. A polynomial of degree $\le2$ in $e^2$ that is not identically zero has at most two zeros. The clause at $\lmu=m$ is \cref{rem:s14:torus-window}.

\emph{(c).} The existence and rotation parts fix $e>0$, $\lambda,m,\lmu>0$ in King's range and the block factor $5$, and fix their constants after the couplings. In their schedule of constants, $\beta$ is fixed last as the minimum of finitely many caps, each one-sided and monotone in $\beta$ with a right side that is a positive finite function of the couplings and of the entries fixed before it; so an admissible $\beta$ exists at every positive point. Ba{\l}aban's own thresholds and King's $J_{\rm K}$ are imposed on $\beta$ after $M$ is fixed.

\emph{(d).} The torus paper states that its numerical window enters through its counterterm-gap input alone and that none of its other conditions is a condition on $e$ or $\lambda$. Its reflection-positivity proposition for King's tori is stated for $\lmu\in(0,\mus]$, $\mus=m$, while the two lemmas on the periodic photon covariance that carry its proof are stated for every $\mu>0$; the proof of \cref{prop:s10:RP} uses only site-independent coefficients, $\lambda>0$ and even sides, and so holds for every $\lmu>0$. The certificate is \cref{rem:s14:torus-window}.

\emph{(e).} The quantifier is read off \cite[Thm.~3.1]{K1} (``for some fixed integer $J$, for all $K\ge J$, and all $k\le K-J$'') and the stopping-scale caps listed above; the last sentence is \cref{rem:s14:quantifier}.
\end{proof}

\begin{remark}[the mass clause of Ba{\l}aban's upper bound]\label{rem:s14:BALm}
Proposition~3.1 of \cite{Ba2} asserts an inequality (3.26), summed over the scales $k=0,\dots,K$, with a constant $\gamma_0=\gamma_0(d,a)$ and mass terms $(L^k\eps)^dm^2|\phi_k|^2$; rescaled to the unit lattice, its $k$-th term is the per-scale inequality (3.29) with the mass term $(L^k\eps)^2m^2|\phi'_k|^2$, so $(L^k\eps)^2m^2$ is the dimensionless mass at scale $k$. The proposition carries the sentence ``we assume $m^2\le O(1)$ also'', $m^2$ being the bare scalar mass squared in Ba{\l}aban's units ($\lambda=1$), i.e.\ $m^2/\lambda^2$ in ours. The clause under which this paper cites the proposition concerns the scales at which it is applied:
\begin{quote}
\emph{Inequality (3.29) of \cite{Ba2} holds, with $\gamma_0=\gamma_0(d,a)$, at every scale $k\le K$ at which $(L^k\eps)^2m^2\le c_m(d,a)$; hence (3.26) holds whenever this is so for every $k\le K$.}
\end{quote}
Here $c_m(d,a)$ is a constant depending on $d$ and $a$ only. The clause belongs to the statement of the proposition: the mass that enters its printed proof, referred to \cite[Sect.~4]{Ba4}, is the unit-lattice mass at scale $k$, the text of \cite{Ba2} makes no further use of the sentence, and the proposition is cited here under the clause. It is then one more one-sided cap on the stopping scale $\beta=\eps_0$ ($(L^k\eps)^2m^2\le\beta^2m^2$ at every scale used), of the kind every other smallness condition of \cite{Ba1,Ba2,Ba3,Ba4} has; a further monotone cap on $\beta$ leaves the existence of an admissible stopping scale unchanged, and nothing is imposed on the bare mass.

Read instead as an absolute bound on $m^2/\lambda^2$, the sentence would exclude small $\lambda$ at fixed $m$, including all of the window $\lambda<0.1m$ of \cite{SymPhase} ($m^2/\lambda^2>100$), and would be an unprinted hypothesis of the theorem (1.3) of \cite{Ba2}, whose only printed condition on $m$ is $m^2>0$; that is not the reading under which the proposition is proved, and it is used nowhere. The clause sits upstream of \cref{hyp:s03:T1} at both masses, of \cref{hyp:s03:EC,hyp:s03:W1} (through \cite[Thm.~3.1]{K1}) and of the steps of the rotation part that rest on \cite{Ba1,Ba2}, as every cap on the stopping scale does; it is not a condition of this paper.
\end{remark}

\begin{remark}[the torus paper's window and its certificate at $\lmu=m$]\label{rem:s14:torus-window}
(i) \emph{The certificate.} The torus paper's coefficient $a_{K,T}(M)=M^2+\dm_{K,T}(\eps;M)$ is the counterterm of \cite[(3.32)--(3.33)]{K1} (orders $2\le\alpha+2\beta\le4$, zero-momentum normalisation) computed on the periodic torus with the periodic scalar propagator of mass $M$ and the periodic photon propagator of the reference mass $\mus$, without block averaging; its theorem on this coefficient gives, for $m>0$, $\mus=m$, $M_1=0.7m$, uniformly in $\eps\le1/m$ and in every periodic torus of side $\ge96/m$, the explicit two-loop bound $a_{K,T}(m)-a_{K,T}(0.7m)\ge0.319\,m^2$. King's torus-convention coefficient $\alpha^{\mathrm{tor}}_{k,T}(M)$ of \cref{def:s10:torus} is the same polynomial over the same orders, as translation-invariant lattice graphs on the torus with the photon of mass $\lmu$ and the same normalisation, the counterterm graphs of \cite{K1} being defined with the full lattice propagators, without block averaging, as we read that paper. So at $\lmu=m=\mus$ the two schemes compute the same graphs and $a_{K,T}(M)=\alpha^{\mathrm{tor}}_{k,T}(M)$ (an identification of ours). Then $s_{k,T}(0.7m)=\frac12(a_{K,T}(m)-a_{K,T}(0.7m))\ge0.1595\,m^2$ for all $\eps_k\le1/m$ and every $T_J$ with $2MJ\ge96/m$; \cref{lem:s10:V} ($J\to\infty$ at fixed $k$) gives $s^\infty_k(0.7m)\ge0.1595\,m^2$ for every such $k$, hence $\sigma_1(0.7m)\ge0.1595\,m^2>0$: at $\lmu=m$, \Tthree\ holds on the whole window with $M_1=0.7m$ and an explicit margin, without \sxivCT\ and without $C_4$; under \sxivCT\ the window lies outside $E^<_W$.

(ii) \emph{Away from $\lmu=m$.} The torus paper's counterterm is computed at the photon mass $\mus=m$, King's at $\lmu$; the two coefficients differ by a term of order $e^2|m-\lmu|$ (the zero-momentum photon-exchange graph at the two photon masses), so the certificate does not transfer to $\lmu\ne m$. There the comparison is through one-loop order: by \eqref{eq:s14:ec2}, $e_c^2\ge4\pi(m+\lmu)>0.1m$, so the one-loop form of \Tthree\ holds on the torus paper's window at every $\lmu>0$; beyond one loop, \cref{prop:s14:T3} places the window inside the region where \Tthree\ holds for every $\lmu\in[0.03m,m]$, and a box shrinking like $\sqrt{\lmu}$ in $e^2$ for every $\lmu\in(0,m]$; the exact position of $E^<_W$ is not determined.
\end{remark}

\medskip\noindent\textbf{The explicit coupling region.} The sign clause \Tthree\ is proved below on an explicit region of the $(\lambda,e^2)$-plane at every photon mass $\lmu\in(0,m]$, with $N=2$ and the auxiliary mass $M_1=0.7m$. The argument has two parts: the three low orders of King's counterterm, whose contributions to $s^\infty_k(M_1)$ are lattice integrals with exact limits (\cref{lem:s14:struct,lem:s14:low,lem:s14:A,lem:s14:B,lem:s14:C}), and the two orders $e^2\lambda$ and $e^4$, whose contributions are bounded from below, uniformly in $k$, through the two-loop majorants of the torus paper's Appendix~G \cite{SymPhase}, read at the photon mass $\lmu$ in place of its reference mass, with the torus floor removed (\cref{lem:s14:E,lem:s14:TF}), made uniform on intervals of photon masses (\cref{lem:s14:MON}) and, in a weaker closed form, uniform in the photon mass (\cref{lem:s14:CF}). The resulting lower bound $\sigma_{\rm low}$ on $\sigma_1(0.7m)$, \eqref{eq:s14:sigmalow}, is explicit through \cref{tab:s14:T3}.

% s14_t3_low: King's counterterm order by order; the three low orders in closed form (input from sec/s14.tex)
\providecommand{\BZ}{\ensuremath{\mathsf B}}

\subsection*{King's counterterm order by order}

Throughout this subsection and the next, $\eps=\eps_k$, $\Lam=\Lk$, $\BZ=\BZ_\eps\defeq[-\pi/\eps,\pi/\eps)^3$ is the Brillouin zone with $\int_\BZ\defeq\int_\BZ\frac{d^3p}{(2\pi)^3}$, $\hat p_\mu\defeq\frac2\eps\sin\frac{\eps p_\mu}2$ and $\hat p^2\defeq\sum_\mu\hat p_\mu^2$ (periodic in $p$), and for a mass $M>0$
\[
G_M(p)\defeq\frac1{\hat p^2+M^2},\qquad \hat D(p)\defeq\frac1{\hat p^2+\lmu^2},\qquad C_M(x)\defeq\int_\BZ e^{\mathrm ip\cdot x}G_M(p),\qquad D(x)\defeq\int_\BZ e^{\mathrm ip\cdot x}\hat D(p)\quad(x\in\Lam)
\]
are the free scalar and photon propagators of King's action on $\Lam$ (per real scalar component and per photon component; $C_M(0)$ is the free Wick constant at mass $M$, $C_m(0)=\sW_\eps$). $N$ is the number of real components of the scalar field. King's free-boundary counterterm \cite[(3.32)--(3.33)]{K1}, \cite[p.~330]{K2} at scalar mass $M$ is the polynomial
\begin{equation}\label{eq:s14:CTdef}
\dm_\eps(M)=\sum_{(\alpha,\beta)}e^\alpha\lambda^\beta\,\dm_{(\alpha\beta)}(\eps;M),\qquad(\alpha,\beta)\in\{(0,1),(2,0),(0,2),(2,1),(4,0)\},
\end{equation}
over the superficially divergent orders $2\le\alpha+2\beta\le4$ with the odd powers of $e$ dropped (they vanish: an odd power of $e$ comes with an odd number of photon fields, whose centred Gaussian expectation is zero); the coefficient $\dm_{(\alpha\beta)}(\eps;M)$ is fixed recursively by the condition that the order-$e^\alpha\lambda^\beta$ part of the amputated one-particle-irreducible two-point function of $\phi$ at zero external momentum, computed on $\Lam$ with the propagators above and with all counterterm vertices of lower order inserted, be cancelled by the counterterm vertex of that order. In the sign convention in which a term of the expansion of $e^{-V}$ that reduces to $+\frac12\sum_{x,y}\eps^6\phi_a(x)\mathsf K(x-y)\phi_a(y)$ contributes $+\hat{\mathsf K}(p)$ to the self-energy $\Sigma(p)$, so that the inverse propagator per component is $\hat p^2+M^2-\Sigma(p)$ and the counterterm vertex $-\frac12\dm\sum\eps^3|\phi|^2$ contributes $-\dm$, the condition reads $\dm_{(\alpha\beta)}(\eps;M)=\Sigma^{(\alpha\beta)}(0;M)$, $\Sigma^{(\alpha\beta)}$ the order-$e^\alpha\lambda^\beta$ part of the sum of the one-particle-irreducible two-point graphs built from the vertices of King's action and from the counterterm vertices of lower order. Every such graph is a finite sum of absolutely convergent integrals over $\BZ^L$, $L\le2$ (all masses are $\ge\min(M,\lmu)>0$ and $\BZ$ is compact), so each coefficient is a finite real number, smooth in $M$ (a rational function of $M^2$ without poles on $(0,\infty)$ under each integral). With $\alpha^\infty_k(M)=M^2+\dm_{\eps_k}(M)$ as in \cref{hyp:s03:T2} and $M_1\in(0,m)$:

\begin{lemma}[structure]\label{lem:s14:struct}
Put $S_{\alpha\beta,k}\defeq\frac12\big[\dm_{(\alpha\beta)}(\eps_k;m)-\dm_{(\alpha\beta)}(\eps_k;M_1)\big]$, numbers free of the couplings. Then
\begin{equation}\label{eq:s14:struct}
s^\infty_k(M_1)=\tfrac12(m^2-M_1^2)+\lambda S_{01,k}+e^2S_{20,k}+\lambda^2S_{02,k}+e^2\lambda S_{21,k}+e^4S_{40,k}.
\end{equation}
In $d=3$, $S_{01,k}$ and $S_{20,k}$ have the dimension of a mass and $S_{02,k}$, $S_{21,k}$, $S_{40,k}$ are dimensionless.
\end{lemma}

\begin{proof}
$s^\infty_k(M_1)=\frac12(\alpha^\infty_k(m)-\alpha^\infty_k(M_1))$ and \eqref{eq:s14:CTdef}.
\end{proof}

\begin{lemma}[the three low orders, exact]\label{lem:s14:low}
For every $\eps$ and $M>0$:
\begin{enumerate}[label=\textup{(\roman*)},leftmargin=*,nosep]
\item $\dm_{(01)}(\eps;M)=-4(N+2)\,C_M(0)$;
\item $\dm_{(20)}(\eps;M)=\int_\BZ\hat k^2\hat D(k)G_M(k)\,dk-3D(0)=-2D(0)-M^2\int_\BZ\hat DG_M$;
\item $\dm_{(02)}(\eps;M)=32(N+2)\sum_{x\in\Lam}\eps^3C_M(x)^3$.
\end{enumerate}
Consequently, with $J_\eps\defeq\int_\BZ G_mG_{M_1}$, $I_\eps\defeq\int_\BZ\hat k^2\hat D(k)G_m(k)G_{M_1}(k)\,dk$ and $F_\eps\defeq\sum_{x\in\Lam}\eps^3[C_{M_1}(x)^3-C_m(x)^3]$,
\begin{equation}\label{eq:s14:lowcoef}
S_{01,k}=2(N+2)\big(C_{M_1}(0)-C_m(0)\big)=2(N+2)(m^2-M_1^2)J_{\eps_k},\qquad S_{20,k}=-\tfrac12(m^2-M_1^2)I_{\eps_k},\qquad S_{02,k}=-16(N+2)F_{\eps_k}.
\end{equation}
\end{lemma}

\begin{proof}
Write $\phi=(\phi_1,\dots,\phi_N)$ with $\phi\cdot\psi=\sum_a\phi_a\psi_a$ and, for $N=2$, $\phi\wedge\psi=\phi_1\psi_2-\phi_2\psi_1=J_{ab}\phi_a\psi_b$ with $J=\begin{pmatrix}0&1\\-1&0\end{pmatrix}$, $J^2=-1$. With $y=x+\eps e_\mu$ and $\theta=e\eps A_\mu(x)$ the hopping term of King's action is $\frac12\eps|e^{\mathrm i\theta}\phi(y)-\phi(x)|^2=\frac12\eps(|\phi(x)|^2+|\phi(y)|^2)-\eps\cos\theta\,\phi(x)\cdot\phi(y)+\eps\sin\theta\,\phi(x)\wedge\phi(y)$; expanding $\cos\theta$ and $\sin\theta$ gives the free action $S_0=\frac12\sum_x\eps^3[\sum_\mu|\nabla_\mu\phi|^2+M^2|\phi|^2]+\frac12\sum_x\eps^3[\sum_{\mu,\nu}(\nabla_\nu A_\mu)^2+\lmu^2|A|^2]$, $\nabla_\mu f(x)=(f(x+\eps e_\mu)-f(x))/\eps$, whose covariances are $\langle\phi_a(x)\phi_b(y)\rangle_0=\delta_{ab}C_M(x-y)$ and $\langle A_\mu(x)A_\nu(y)\rangle_0=\delta_{\mu\nu}D(x-y)$, and the vertices
\[
V_1=e\sum_{x,\mu}\eps^2A_\mu(x)\,\phi(x)\wedge\phi(y),\qquad V_2=\tfrac{e^2}2\sum_{x,\mu}\eps^3A_\mu(x)^2\,\phi(x)\cdot\phi(y),\qquad V_\lambda=\lambda\sum_x\eps^3(\phi\cdot\phi)^2,\qquad V_{\rm ct}=\tfrac12\dm\sum_x\eps^3\phi\cdot\phi,
\]
together with vertices of order $e^3$, $e^4$ and higher, which do not enter the three orders below. The graphs are those of $-V+\frac12V^2-\cdots$ with the normalisation fixed above; by $O(N)$ symmetry they are proportional to $\delta_{ab}$.

(i) At order $\lambda$ the only one-particle-irreducible two-point graph is the tadpole of $-V_\lambda$. In $(\phi\cdot\phi)^2=\phi_a\phi_a\phi_b\phi_b$ one of the six pairs of fields is contracted: the pairs $(1,2)$ and $(3,4)$ each give $N\,C_M(0)\,\phi\cdot\phi$, the four mixed pairs each give $C_M(0)\,\phi\cdot\phi$, in all $2(N+2)C_M(0)\,\phi\cdot\phi$. So $-V_\lambda$ contributes $-2(N+2)\lambda C_M(0)\sum_x\eps^3\phi\cdot\phi=\frac12\sum\phi_a\mathsf K\phi_a$ with $\hat{\mathsf K}=-4(N+2)\lambda C_M(0)$, and $\dm_{(01)}=\Sigma^{(01)}(0)/\lambda=-4(N+2)C_M(0)$.

(ii) At order $e^2$ there are two graphs. \emph{Seagull:} contracting $A_\mu(x)^2$ to $D(0)$ in $-V_2$ gives $-\frac{e^2}2D(0)\sum_{x,\mu}\eps^3\phi(x)\cdot\phi(x+\eps e_\mu)$, i.e.\ $\Sigma_{\rm seag}(p)=-e^2D(0)\sum_\mu\cos(\eps p_\mu)$ after symmetrisation, and $\Sigma_{\rm seag}(0)=-3e^2D(0)$. \emph{Exchange:} in momentum space $V_1=\sum_\mu\iiint\hat A_\mu(k)\hat\phi_i(p)\hat\phi_j(p')\,W^\mu_{ij}(p,p')\,(2\pi)^3\delta(k+p+p')$ with the symmetrised vertex factor $W^\mu_{ij}(p,p')=\mathrm ie\,J_{ij}\,e^{-\mathrm i\eps k_\mu/2}\frac1\eps\sin\frac{\eps(p'-p)_\mu}2$ (from $\frac12J_{ij}(e^{\mathrm i\eps p'_\mu}-e^{\mathrm i\eps p_\mu})=J_{ij}e^{\mathrm i\eps(p+p')_\mu/2}\,\mathrm i\sin\frac{\eps(p'-p)_\mu}2$). In $\frac12V_1^2$ contract the photon pair, $\langle\hat A_\mu\hat A_\nu\rangle=\delta_{\mu\nu}(2\pi)^3\delta(k+k')\hat D(k)$, and one scalar pair (four equivalent choices); with external pair $\hat\phi_i(p)\hat\phi_l(-p)$ the internal line carries $p+k$ and the two vertex factors multiply to $(\mathrm ie)^2(J^2)_{il}\frac1{\eps^2}\sin^2\frac{\eps(2p+k)_\mu}2=e^2\delta_{il}\frac1{\eps^2}\sin^2\frac{\eps(2p+k)_\mu}2$, the phases cancelling. Hence $\Sigma_{\rm exch}(p)=e^2\int_\BZ\hat D(k)G_M(p+k)\sum_\mu\big(\frac2\eps\sin\frac{\eps(2p+k)_\mu}2\big)^2dk$ and $\Sigma_{\rm exch}(0)=e^2\int_\BZ\hat D(k)\hat k^2G_M(k)\,dk$. So $\dm_{(20)}=\int_\BZ\hat k^2\hat DG_M-3D(0)$; the second form follows from $\hat k^2\hat DG_M=\hat D(1-M^2G_M)$ and $\int_\BZ\hat D=D(0)$.

(iii) At order $\lambda^2$ (and $e=0$) the one-particle-irreducible two-point graphs with two quartic vertices $u,v$ are the sunset (external legs at $u$ and at $v$, three lines $u$--$v$) and the tadpole with a tadpole inserted on its loop line (external legs at $u$, two lines $u$--$v$, one loop at $v$); the configurations with a single bridge are not one-particle irreducible and the rest are disconnected. The insertion graph is cancelled exactly by the graph with the counterterm vertex $-V^{(01)}_{\rm ct}$ on the loop line, because the order-$\lambda$ self-energy with its counterterm, $\Sigma^{(01)}(q)-\lambda\dm_{(01)}$, vanishes for every $q$ (the tadpole does not depend on the momentum). For the sunset, in $\frac12V_\lambda^2=\frac12\lambda^2\sum_{u,v}\eps^6(\phi\cdot\phi)^2_u(\phi\cdot\phi)^2_v$ choose the external field at $u$ ($4$ ways) and at $v$ ($4$ ways) and pair the remaining triples in the $3!$ ways: two pairings give $N\delta_{ab}$ and four give $\delta_{ab}$, in all $2(N+2)\delta_{ab}$; so $\frac12V_\lambda^2$ contributes $\frac12\sum_{u,v}\eps^6\phi_a(u)\mathsf K(u-v)\phi_a(v)$ with $\mathsf K(x)=32(N+2)\lambda^2C_M(x)^3$, and $\dm_{(02)}=\hat{\mathsf K}(0)/\lambda^2=32(N+2)\sum_x\eps^3C_M(x)^3$.

The formulas \eqref{eq:s14:lowcoef} follow from (i)--(iii) and the resolvent identity $G_{M_1}-G_m=(m^2-M_1^2)G_{M_1}G_m$; the $M$-independent tadpole $3D(0)$ and the term $-2D(0)$ cancel in $S_{20,k}$.
\end{proof}

At $N=2$ the three coefficients are the terms $-16\lambda C_T(M)$, $-2e^2D_T-e^2J_T(M)$ and $128\lambda^2S_T(M)$ of the exact two-loop form of the torus coefficient in the torus paper \cite[App.~G, Lemma~G.1]{SymPhase}, read on $\Lam$ in place of the torus; the orders $e^2\lambda$ and $e^4$ are the terms $16e^2\lambda F_T(M)$ and $e^4Q_T(M)$ there. Every graph of these two orders is an absolutely convergent integral over $\BZ^2$ at fixed $\eps$ (massive propagators, bounded trigonometric vertex factors), so $S_{21,k}$ and $S_{40,k}$ are finite real numbers for every $k$; their bounds, uniform in $k$, are the subject of the next subsection.

% s14_t3_low2: the one-loop coefficients and the sunset; limits and k-uniform bounds (input from sec/s14.tex)

\begin{lemma}[four lattice facts]\label{lem:s14:L}
\textup{(L1)} For $p\in\BZ$: $c\,p^2\le\hat p^2\le p^2$ with $c=4/\pi^2$. Hence, for a mass $a>0$, $\bar G_a(p)\defeq(p^2+a^2)^{-1}\le G_a(p)\le\bar G^c_a(p)\defeq(cp^2+a^2)^{-1}$ on $\BZ$, and for $a,b>0$: $\int_\BZ\bar G^c_a\bar G^c_b\le\int_{\R^3}\bar G^c_a\bar G^c_b\,\frac{d^3p}{(2\pi)^3}=\frac{\pi^2}{32(a+b)}$, $\int_\BZ(\bar G^c_a)^2\le\frac{\pi^2}{64a}$.
\textup{(L2)} $C_M(x)=\eps^{-3}\int_0^\infty e^{-tM^2}(e^{t\Delta_\eps})_{x,0}\,dt$ with $(e^{t\Delta_\eps})_{x,y}=e^{-6t/\eps^2}\sum_{n\ge0}\frac{(t/\eps^2)^n}{n!}\#\{n\text{-step nearest-neighbour paths }y\to x\}>0$. Hence $C_M(x)>0$ for every $x\in\Lam$, $M\mapsto C_M(x)$ is strictly decreasing, $\sum_x\eps^3C_M(x)=M^{-2}$, and the resolvent identity $C_{M_1}-C_m=(m^2-M_1^2)\,C_{M_1}*C_m$ holds pointwise with a positive right side, $(f*g)(x)\defeq\sum_y\eps^3f(x-y)g(y)$.
\textup{(L3)} With $\int\defeq\int_{\R^3}\frac{d^3p}{(2\pi)^3}$ and $a,b,c>0$: $\int\frac1{(p^2+a^2)(p^2+b^2)}=\frac1{4\pi(a+b)}$, $\int\frac1{(p^2+a^2)^2}=\frac1{8\pi a}$, $\int\frac1{(p^2+a^2)(p^2+b^2)(p^2+c^2)}=\frac1{4\pi(a+b)(b+c)(c+a)}$.
\textup{(L4)} If $h_\eps$ on $\BZ_\eps$ satisfies $|h_\eps|\le\bar h\in L^1(\R^3)$ and $h_\eps(p)\to h(p)$ for every $p$ as $\eps\to0$, then $\int_{\BZ_\eps}h_\eps\to\int_{\R^3}h$. As $\eps\to0$, $\hat p^2\to p^2$ at every fixed $p$, so $G_M(p)\to\bar G_M(p)$ and $\hat D(p)\to\bar G_{\lmu}(p)$.
\end{lemma}

\begin{proof}
(L1) $|\sin x|\le|x|$ and $\sin x\ge2x/\pi$ on $[0,\pi/2]$, at $x=\eps|p_\mu|/2\le\pi/2$; the integrals follow from (L3) by the scaling $p\mapsto c^{-1/2}p$, $c^{-3/2}=\pi^3/8$. (L2) $(-\Delta_\eps+M^2)^{-1}=\int_0^\infty e^{-tM^2}e^{t\Delta_\eps}dt$ as matrices on $\Lam$, with kernel $\eps^3C_M(x-y)$; the path expansion is $e^{t\Delta_\eps}=e^{-6t/\eps^2}e^{(t/\eps^2)\mathcal A}$ with $\mathcal A$ the adjacency matrix; $\sum_x(e^{t\Delta_\eps})_{x,0}=1$; the identity is $G_{M_1}-G_m=(m^2-M_1^2)G_{M_1}G_m$ in Fourier space. (L3) Radially, $\frac1{2\pi^2}\int_0^\infty\frac{p^2dp}{(p^2+a^2)(p^2+b^2)}=\frac1{2\pi^2}\cdot\frac{\pi}{2(a+b)}$, the second integral is the limit $b\to a$, and the third follows by partial fractions from the first. (L4) Dominated convergence after extension by $0$.
\end{proof}

\begin{lemma}[order $\lambda$]\label{lem:s14:A}
For every $k$: \textup{(i)} $S_{01,k}>0$; \textup{(ii)} $J_\eps\le\frac{\pi^2}{32(m+M_1)}$, hence $S_{01,k}\le\frac{(N+2)\pi^2(m-M_1)}{16}$; \textup{(iii)} $J_{\eps_k}\to\frac1{4\pi(m+M_1)}$, hence $S_{01,k}\to\frac{(N+2)(m-M_1)}{2\pi}$ as $k\to\infty$.
\end{lemma}

\begin{proof}
(i) \eqref{eq:s14:lowcoef} and $G_mG_{M_1}>0$. (ii) (L1). (iii) (L4) with $\bar h=\bar G^c_m\bar G^c_{M_1}$, and (L3).
\end{proof}

\begin{lemma}[order $e^2$]\label{lem:s14:B}
For every $k$: \textup{(i)} $S_{20,k}<0$; \textup{(ii)} $0<I_\eps\le\bar I\defeq\frac{\pi^2}{32(m+M_1)}$, a bound free of $\lmu$ and $k$; \textup{(iii)} $I_{\eps_k}\to I_\infty(\lmu)\defeq\frac{m\lmu+mM_1+M_1\lmu}{4\pi(m+M_1)(\lmu+M_1)(m+\lmu)}$, hence $S_{20,k}\to-\frac12(m^2-M_1^2)I_\infty(\lmu)=-\frac{(m-M_1)(m\lmu+mM_1+M_1\lmu)}{8\pi(\lmu+M_1)(m+\lmu)}$; \textup{(iv)} $\lmu\mapsto I_\infty(\lmu)$ is strictly decreasing on $(0,\infty)$, with $I_\infty(\lmu)\uparrow I^*\defeq\frac1{4\pi(m+M_1)}$ as $\lmu\downarrow0$.
\end{lemma}

\begin{proof}
(i) \eqref{eq:s14:lowcoef}. (ii) $\hat k^2\hat D(k)\in[0,1)$ and (L1). (iii) (L4) with $\bar h=\bar G^c_m\bar G^c_{M_1}$; the limit is $\int\frac{k^2}{(k^2+\lmu^2)(k^2+m^2)(k^2+M_1^2)}=\frac1{4\pi(m+M_1)}-\frac{\lmu^2}{4\pi(\lmu+m)(m+M_1)(M_1+\lmu)}$ by (L3), and $(\lmu+M_1)(m+\lmu)-\lmu^2=m\lmu+mM_1+M_1\lmu$. (iv) With $\mathcal N(\mu)=(m+M_1)\mu+mM_1$ and $\mathcal D(\mu)=(\mu+M_1)(m+\mu)$, $\mathcal N'\mathcal D-\mathcal N\mathcal D'=-(m+M_1)\mu^2-2mM_1\mu<0$.
\end{proof}

\begin{lemma}[order $\lambda^2$]\label{lem:s14:C}
For every $k$: \textup{(i)} $F_\eps>0$, so $S_{02,k}<0$; \textup{(ii)} $F_\eps\le3(m^2-M_1^2)\big(\int_\BZ G_{M_1}^2\big)^2\le\frac{3\pi^4(m^2-M_1^2)}{4096M_1^2}$, hence $|S_{02,k}|\le\bar L\defeq\frac{3(N+2)\pi^4(m^2-M_1^2)}{256M_1^2}$; \textup{(iii)} $F_{\eps_k}\to\frac{\ln(m/M_1)}{16\pi^2}$, hence $S_{02,k}\to-L(M_1)\defeq-\frac{(N+2)\ln(m/M_1)}{\pi^2}$. The same holds with any $M>0$ in place of $M_1$: $\lim_k[\dm_{(02)}(\eps_k;m)-\dm_{(02)}(\eps_k;M)]=-\frac{N+2}{\pi^2}\ln\frac{m^2}{M^2}$ (each $\dm_{(02)}(\eps_k;M)$ alone diverges like $\ln(1/\eps_k)$).
\end{lemma}

\begin{proof}
(i) By (L2), $C_{M_1}(x)>C_m(x)>0$ pointwise, and the sum converges absolutely ($C_M\in\ell^1\cap\ell^\infty$).

(ii) Pointwise $C_{M_1}^3-C_m^3=(C_{M_1}-C_m)(C_{M_1}^2+C_{M_1}C_m+C_m^2)\le3(C_{M_1}-C_m)C_{M_1}^2$, and by (L2) $C_{M_1}-C_m=(m^2-M_1^2)C_{M_1}*C_m\le(m^2-M_1^2)C_{M_1}*C_{M_1}$. For $f,g,h\in\ell^1(\Lam)$, $\sum_x\eps^3f(x)g(x)h(x)=\int_\BZ\!\int_\BZ\hat f(p)\hat g(q)\hat h(-p-q)$ (Fourier inversion and Fubini, $\hat h$ periodic); here $\hat f=G_{M_1}^2$, $\hat g=\hat h=G_{M_1}$, all even. By $G(q)G(p+q)\le\frac12[G(q)^2+G(p+q)^2]$ and the measure-preserving shift $q\mapsto p+q$ of $\BZ$,
\[
\int_\BZ\!\int_\BZ G_{M_1}(p)^2G_{M_1}(q)G_{M_1}(p+q)\le\Big(\int_\BZ G_{M_1}^2\Big)^2\le\Big(\frac{\pi^2}{64M_1}\Big)^2
\]
by (L1). Multiply by $3(m^2-M_1^2)$, and by $16(N+2)$ for $\bar L$.

(iii) Write $G_1\defeq G_{M_1}$, $G_2\defeq G_m$ (at spacing $\eps$) and $\Delta\defeq G_1-G_2=(m^2-M_1^2)G_1G_2\in[0,(m^2-M_1^2)G_1^2]$, so that, by the formula of (ii), $F_\eps=\int_\BZ\!\int_\BZ[G_1G_1G_1-G_2G_2G_2](p,q,p+q)=T_1+T_2+T_3$ with $T_1=\iint\Delta(p)G_1(q)G_1(p+q)$, $T_2=\iint G_2(p)\Delta(q)G_1(p+q)$, $T_3=\iint G_2(p)G_2(q)\Delta(p+q)$, all $\ge0$ (telescoping). Consider $T_3$; the other two are the same argument after the measure-preserving change of variables $(p,q)\mapsto(q,-p-q\bmod\BZ)$ on $\BZ\times\BZ$. Split $\BZ\times\BZ=N_\eps\cup U_\eps$ with $U_\eps\defeq\{(p,q):|p_\mu+q_\mu|>\pi/\eps\text{ for some }\mu\}$, the set where $p+q$ is folded back into $\BZ$. \emph{The folding set is negligible:} on $U_\eps$, $\max(|p|,|q|)\ge\pi/(2\eps)$, so by (L1) $\min(G_2(p),G_2(q))\le(c\pi^2/(4\eps^2))^{-1}=\eps^2$; with $\Delta\le(m^2-M_1^2)G_1G_2$, the substitution $r=p+q\pmod\BZ$ at fixed $q$, and $\int_\BZ G_2\le\int_{|p|\le\sqrt3\pi/\eps}\frac{d^3p}{(2\pi)^3}\frac1{cp^2}=\frac{\sqrt3\pi}{8\eps}$ (L1),
\[
\iint_{U_\eps}G_2(p)G_2(q)\Delta(p+q)\le2(m^2-M_1^2)\,\eps^2\int_\BZ G_2\int_\BZ G_1G_2\le2(m^2-M_1^2)\,\eps^2\cdot\frac{\sqrt3\pi}{8\eps}\cdot\frac{\pi^2}{32(m+M_1)}\xrightarrow[\eps\to0]{}0 .
\]
\emph{The rest converges:} on $N_\eps$ the integrand, extended by $0$, converges pointwise (every fixed $(p,q)$ lies in $N_\eps$ for small $\eps$, and (L4)) to $(m^2-M_1^2)\bar G_m(p)\bar G_m(q)\bar G_{M_1}(p+q)\bar G_m(p+q)$ and is dominated by $(m^2-M_1^2)\bar G^c_{M_1}(p)\bar G^c_{M_1}(q)\bar G^c_{M_1}(p+q)^2\in L^1(\R^6)$ (substitute $r=p+q$ and use $\bar G^c(q)\bar G^c(r-q)\le\frac12[\bar G^c(q)^2+\bar G^c(r-q)^2]$; the integral is at most $(\int(\bar G^c_{M_1})^2)^2<\infty$). So $T_3\to T_3^\infty\defeq(m^2-M_1^2)\iint_{\R^6}\bar G_m(p)\bar G_m(q)\bar G_{M_1}(p+q)\bar G_m(p+q)$, and likewise $T_1$, $T_2$. \emph{Back to position space:} with the continuum Green's function $\bar C_M(x)=\frac{e^{-M|x|}}{4\pi|x|}$ (Fourier transform $\bar G_M$) and $\hat h\defeq\bar G_{M_1}\bar G_m\in L^1\cap L^2$, $h=\bar C_{M_1}*\bar C_m=(\bar C_{M_1}-\bar C_m)/(m^2-M_1^2)$, a bounded continuous function in $L^2$: for fixed $p$, $\int\bar G_m(q)\hat h(p+q)\,dq=\widehat{(\bar C_mh)}(p)$ as a Plancherel identity for the $L^2$ functions $\bar G_m$ and $\hat h$ ($\bar C_mh\in L^1\cap L^2$, being $O(|x|^{-1})$ near $0$ and exponentially small at infinity), and $\int\bar G_m(p)\widehat{(\bar C_mh)}(p)\,dp=\int d^3x\,\bar C_m(x)(\bar C_mh)(x)$ by Plancherel again ($\bar C_m\in L^2(\R^3)$). Hence $T_3^\infty=\int d^3x\,\bar C_m^2(\bar C_{M_1}-\bar C_m)$, $T_1^\infty=\int\bar C_{M_1}^2(\bar C_{M_1}-\bar C_m)$, $T_2^\infty=\int\bar C_m\bar C_{M_1}(\bar C_{M_1}-\bar C_m)$, each absolutely convergent (the integrands are $O(|x|^{-2})$ at $0$), and
\[
\lim_{\eps\to0}F_\eps=\int_{\R^3}d^3x\,\big[\bar C_{M_1}(x)^3-\bar C_m(x)^3\big]=\frac{4\pi}{(4\pi)^3}\int_0^\infty\frac{e^{-3M_1x}-e^{-3mx}}{x}\,dx=\frac{\ln(m/M_1)}{16\pi^2}
\]
by Frullani's integral $\int_0^\infty(e^{-ax}-e^{-bx})\frac{dx}x=\ln\frac ba$. Multiplying by $-16(N+2)$ gives the limit of $S_{02,k}$. For an arbitrary $M>0$ in place of $M_1$ the same proof applies (for $M>m$ with the roles of the two masses exchanged); the divergence of a single $\dm_{(02)}(\eps;M)$ is that of $\sum_x\eps^3C_M(x)^3$, whose terms near $x=0$ are of size $\eps^{-3}$ on $O(1)$ sites, with the logarithm from the region $\eps\ll|x|\ll1/M$, where $C_M(x)\approx(4\pi|x|)^{-1}$.
\end{proof}

\begin{remark}[the one-loop form and the $e=0$ limit]\label{rem:s14:lowlimits}
\Cref{lem:s14:A,lem:s14:B,lem:s14:C} give the limits $\lim_kS_{01,k}=\frac{(N+2)(m-M_1)}{2\pi}$, $\lim_kS_{20,k}=-\frac12(m^2-M_1^2)I_\infty(\lmu)$ and $\lim_kS_{02,k}=-L(M_1)$ for every $M_1\in(0,m)$ and every $\lmu>0$: inserted into \eqref{eq:s14:struct} they are the one-loop closed form of \cref{lem:s14:CW}(b)(ii) and, at $e=0$, the exact limit $\lim_ks^\infty_k(M_1)=\frac12(m^2-M_1^2)+\frac{(N+2)\lambda(m-M_1)}{2\pi}-\frac{(N+2)\lambda^2}{\pi^2}\ln\frac m{M_1}$ of \cref{prop:s14:T3}(vi). The constants $\bar I$, $\bar L$ and the limits of $S_{01,k}$, $S_{02,k}$ do not involve the photon mass; of the three low orders only $I_\infty$ depends on $\lmu$, and $I_\infty(\lmu)\le I^*$ for every $\lmu>0$ (\cref{lem:s14:B}(iv)).
\end{remark}

% s14_t3_maj: the two-loop majorants at every photon mass: torus floor, monotonicity, closed form (input from sec/s14.tex)

\subsection*{The two-loop orders at every photon mass}

From here to the end of the proof of \cref{prop:s14:T3} we use the units $m=1$ ($\lambda$ and $e^2$ in units of $m$, $\sigma$ in units of $m^2$), $N=2$ and $M_1=0.7$; ``$\eps_k\le1$'' means $\eps_k\le1/m$. The two-loop coefficients $S_{21,k}$, $S_{40,k}$ are bounded from below through the torus paper's Appendix~G \cite{SymPhase}, whose statements are applied with their hypotheses matched here; that appendix is written on periodic tori $T$ of side $\ell_T\ge\ell_0$ with $\ell_T/\eps\in2\N$, in units $m=1$, and its statements carry the reference photon mass $\mu_*\in(0,1]$ of its photon propagator $\hat D=(\hat k^2+\mu_*^2)^{-1}$ as a parameter; we read them at $\mu_*=\lmu$. Its objects are the following.

\begin{enumerate}[label=\textup{(G\arabic*)},leftmargin=*]
\item \emph{The torus coefficient} $a_{K,T}(M)$ \cite[Def.~2.2]{SymPhase}: King's recursion \eqref{eq:s14:CTdef} computed on $T$ with the periodic propagators of $T$ at scalar mass $M$ and photon mass $\mu_*$. Read at $\mu_*=\lmu$, this is the bare coefficient $\alpha^{\mathrm{tor}}_{k,T}(M)=M^2+\dm_{k,T}(M)$ of King's torus convention (\cref{def:s10:torus}(ii)): both are the same recursion with the same propagators. Its exact two-loop form \cite[Lemma~G.1]{SymPhase} is, at $N=2$,
\[
a_{K,T}(M)=M^2-16\lambda C_T(M)-2e^2D_T-e^2J_T(M)+128\lambda^2S_T(M)+16e^2\lambda F_T(M)+e^4Q_T(M),
\]
with $C_T(M)=\int_T\hat G_M$, $D_T=\int_T\hat D$, $J_T(M)=M^2\int_T\hat D\hat G_M$, $S_T(M)=\sum_{x\in T}\eps^3G_M(x)^3$, $F_T$ the order-$e^2\lambda$ and $Q_T$ the order-$e^4$ coefficient, the term $\frac34\eps^2D_T^2$ of $Q_T$ being independent of $M$; $M\mapsto a_{K,T}(M)$ is smooth. (The three low orders are \cref{lem:s14:low} on the torus.)
\item \emph{The one-loop majorants} \cite[Lemma~G.15]{SymPhase}: for $\alpha,\beta\ge0$ with $g\defeq\alpha+\beta>\frac32$, $a,b>0$, $\nu\defeq\min(a,b)$ and $\ell\in(0,\infty]$,
\begin{equation}\label{eq:s14:Iab}
\mathcal I_{\alpha,\beta}(a,b;\ell)\defeq c_NC_3+\frac{3\pi}{16\ell}C_2+\frac3{4\ell^2}C_1+\frac{a^{-2\alpha}b^{-2\beta}}{\ell^3},\qquad c_N=\frac{\pi^3}8,
\end{equation}
each $C_d=C_d(\alpha,\beta;a,b)$ the minimum of its admissible closed forms: for $C_3$ the $\nu$-form $\Gamma(g-\frac32)\nu^{3-2g}/((4\pi)^{3/2}\Gamma(g))$, the Beta form in $a$, $\frac1{4\pi^2}B(c,\alpha-c)a^{2c-2\alpha}$ with $c=\frac{3-2\beta}2$ if $\beta<\frac32$, the same with $(a,\alpha)\leftrightarrow(b,\beta)$ if $\alpha<\frac32$, the exact values $\frac1{4\pi(a+b)}$ at $(\alpha,\beta)=(1,1)$ and $\frac1{8\pi a(a+b)^2}$ at $(2,1)$, and the H\"older form $\gamma_3(g)a^{(3-2g)\alpha/g}b^{(3-2g)\beta/g}$, $\gamma_3(g)a^{3-2g}$ the closed form of $\int(u+a^2)^{-g}$; for $C_2$ the $\nu$-form $\nu^{2-2g}/(g-1)$, the Beta form $a^{2(1-g)}B(1-\beta,g-1)$ if $0<\beta<1$ and its symmetric form, and the H\"older form; for $C_1$ the $\nu$-form $\sqrt\pi\Gamma(g-\frac12)\nu^{1-2g}/\Gamma(g)$, the Beta form $a^{1-2g}B(\frac12-\beta,g-\frac12)$ if $0<\beta<\frac12$ and its symmetric form, and the H\"older form. For every torus of side $\ell$ and every $\eps$, $\int_T(\hat k^2+a^2)^{-\alpha}(\hat k^2+b^2)^{-\beta}\le\mathcal I_{\alpha,\beta}(a,b;\ell)$, on $\eps\Z^3$ with $\ell=\infty$; $\mathcal I_{\alpha,\beta}$ is decreasing in $\ell$, and at $(a,b)=(M,\mu)$ every admissible form is non-increasing in $M$. The majorants, all at $(M,\mu;\ell_0)$ with $\mu=\mu_*$, are $\bar\tau=\mathcal I_{2,0}$, $\bar\psi=\mathcal I_{1,1}$, $\bar\delta=\mathcal I_{0,2}$, $\bar\Theta=\mathcal I_{1,3/4}+\mathcal I_{3/4,1}$, $\bar J_s=\mathcal I_{3-s/2,0}$ and $\bar J'_s=\mathcal I_{2-s/2,1}$ for $s\in\{0,1,\frac32\}$, $\bar\Lambda_0=\mathcal I_{3/4,1}$, $\bar\Lambda_1=\mathcal I_{5/4,1}+\mathcal I_{2,1/4}$, $\bar\Lambda_2=\mathcal I_{1/4,2}+\mathcal I_{1,5/4}$: fourteen distinct pairs $(\alpha,\beta)$, all with $g\ge\frac74>\frac32$; and $\bar\omega\defeq w+\hat E(\mu,\ell_0)$ with $w=0.252731$ (used as $w\le0.25273101$), $\hat E\defeq\min\{\bar E,\tilde E\}$, where $\bar E(\mu,\ell)=\frac{7.37}\ell+\frac{3\pi}{4\ell}E_1(\frac{\mu^2\ell^2}8)+\frac{12\pi\,\mathrm{erfc}(\mu\ell/(2\sqrt2))}{\mu\ell^2}+\frac{64e^{-\mu^2\ell^2/8}}{\mu^2\ell^3}$ \cite[Lemma~H.31]{SymPhase} and $\tilde E\ge0$ \cite[Lemma~G.14]{SymPhase}; for $\eps\le1$ and $\ell_T\ge\ell_0$, $\eps D_T\le\bar\omega$ \cite[Lemma~G.14(c)]{SymPhase}, and on $\eps\Z^3$ ($\ell_0=\infty$) one takes $\bar\omega=w$ \cite[Lemma~G.24]{SymPhase}.
\item \emph{The slope bounds} \cite[Lemma~G.2]{SymPhase}: for $\eps\le1$, $\ell_T\ge\ell_0$, $M>0$, with the majorants at $(M,\mu_*;\ell_0)$, $\bar\beta\defeq2M[2\bar\tau+8\bar\psi]$ and
\begin{gather*}
j\defeq2M\min\{\bar\tau,\bar\psi\},\qquad \varphi_3\defeq4M\big[2\sqrt2\,\bar\Theta\bar J_{3/2}+|\mu_*^2-2M^2|\bar\psi\bar J_0+3\sqrt3\,\bar\omega\bar J_1\big]+\bar\beta\bar\tau,\\
b_4\defeq\varphi_A+28M\bar\delta\bar\tau+\varphi_{D2},\quad \varphi_A\defeq6M\big[2\sqrt2\,\bar\Theta\bar J'_{3/2}+|\mu_*^2-2M^2|\bar\psi\bar J'_0+3\sqrt3\,\bar\omega\bar J'_1\big]+\bar\beta\bar\psi,\\
\varphi_{D2}\defeq2M\bar\psi(2\bar\psi+4\bar\delta)+\min\{8M\bar\psi(6\bar\psi+2\bar\tau+4\bar\delta),\ 2^{5/2}M(2\bar\Lambda_1+4\bar\Lambda_2)\bar\Lambda_0\},
\end{gather*}
one has $0\le\partial_MJ_T\le j(M)$, $0\le-128\,\partial_MS_T\le768M\bar\tau(M)^2$, $|\partial_MF_T|\le\varphi_3(M)$ and $\partial_MQ_T\ge-b_4(M)$.
\item \emph{The two-loop gap} \cite[Prop.~G.3]{SymPhase}: for $0<M_1<M_2$, $\mu_*\in(0,1]$, $\ell_0>0$, every $\eps\le1$ and every torus of side $\ell_T\ge\ell_0$, $a_{K,T}(M_2)-a_{K,T}(M_1)\ge(M_2^2-M_1^2)+16\lambda h_1-e^2\bar j-\lambda^2\bar s-e^2\lambda\bar B_3-e^4\bar B_4^-$ with $h_1\ge0$, $\bar j=\int_{M_1}^{M_2}j$, $\bar s=\int768M\bar\tau^2$, $\bar B_3=16\int\varphi_3$, $\bar B_4^-=\int b_4$; its proof writes each order as $\int_{M_1}^{M_2}\partial_M(\cdot)$ and applies (G3) \emph{order by order}: $-[J_T(M_2)-J_T(M_1)]\ge-\bar j$, $128[S_T(M_2)-S_T(M_1)]\ge-\bar s$, $16[F_T(M_2)-F_T(M_1)]\ge-\bar B_3$, $Q_T(M_2)-Q_T(M_1)\ge-\bar B_4^-$.
\item \emph{The cell rule} (the proof of \cite[Cor.~G.4]{SymPhase}): for a density $\phi\in\{j,768M\bar\tau^2,16\varphi_3,b_4\}$ and a cell $[a,b]\subset(0,\infty)$ let $\phi^{\rm cell}[a,b](\mu,\ell_0)$ be the number obtained from the formula for $\phi$ by evaluating every majorant at the left endpoint $a$, the prefactor $M$ at the right endpoint $b$, the prefactor $|\mu^2-2M^2|$ at its maximum over $\{a,b\}$ and $\bar\omega$ as it is; since the majorants are non-negative and non-increasing in $M$ and the formulas are sums, products and minima with non-negative coefficients, $\phi\le\phi^{\rm cell}[a,b]$ on $[a,b]$, and for a partition $M_1=M^{(0)}<\dots<M^{(n)}=M_2$
\begin{equation}\label{eq:s14:cell}
\int_{M_1}^{M_2}\phi\,dM\ \le\ \Phi_n(\mu,\ell_0)\defeq\sum_{i=0}^{n-1}\big(M^{(i+1)}-M^{(i)}\big)\,\phi^{\rm cell}[M^{(i)},M^{(i+1)}](\mu,\ell_0).
\end{equation}
We write $\bar j_n,\bar s_n,\bar B_{3,n},\bar B^-_{4,n}$ for the four upper sums on the uniform partition of $[0.7,1]$ into $n$ cells, and $\bar X_n$ for any of them; $n=400$ unless said otherwise.
\end{enumerate}

\begin{lemma}[the four orders on $\Lk$ at every photon mass]\label{lem:s14:E}
Let $\lmu\in(0,1]$, $\ell_0>0$, $n\ge1$. For every $k$ with $\eps_k\le1$:
\[
2S_{20,k}\ge-\bar j_n(\lmu,\ell_0),\qquad 2S_{02,k}\ge-\bar s_n,\qquad 2S_{21,k}\ge-\bar B_{3,n}(\lmu,\ell_0),\qquad 2S_{40,k}\ge-\bar B^-_{4,n}(\lmu,\ell_0).
\]
\end{lemma}

\begin{proof}
Fix such a $k$, and let $T=T_J$ be King's tori (\cref{def:s10:torus}; side $2MJ$, $2MJ\cdot5^k$ sites per side, an even number), $J$ so large that $2MJ\ge\ell_0$. By (G1) at $\mu_*=\lmu$, the coefficients of the torus-convention bare coefficient $\alpha^{\mathrm{tor}}_{k,T}(M)$ in the five orders are the terms of the exact two-loop form, so that for $S^T_{\alpha\beta,k}\defeq\frac12[\dm_{(\alpha\beta),T}(\eps_k;1)-\dm_{(\alpha\beta),T}(\eps_k;0.7)]$,
\[
2S^T_{20,k}=-[J_T(1)-J_T(0.7)],\quad 2S^T_{02,k}=128[S_T(1)-S_T(0.7)],\quad 2S^T_{21,k}=16[F_T(1)-F_T(0.7)],\quad 2S^T_{40,k}=Q_T(1)-Q_T(0.7),
\]
the $M$-independent terms $-2D_T$ and $\frac34\eps^2D_T^2$ cancelling in the differences. The hypotheses of (G3) and (G4) are met: $\eps_k\le1$, $\ell_T\ge\ell_0$, $\mu_*=\lmu\in(0,1]$, $M\in[0.7,1]$, and the lemmas behind (G3) carry $\mu_*$ as a parameter in their statements ((G2): $a,b>0$ and $g>\frac32$ for every instance; $\eps\le1$ and $\ell_T\ge\ell_0$ for $\bar\omega$; $\ell_T/\eps\in\N$, $\ell_T/\eps\ge2$ for $\bar E$). Hence, by the order-by-order bounds of (G4) with the majorants at $(\lmu,\ell_0)$ and by the cell rule \eqref{eq:s14:cell}, $2S^T_{20,k}\ge-\bar j_n(\lmu,\ell_0)$, $2S^T_{02,k}\ge-\bar s_n$, $2S^T_{21,k}\ge-\bar B_{3,n}(\lmu,\ell_0)$, $2S^T_{40,k}\ge-\bar B^-_{4,n}(\lmu,\ell_0)$, uniformly in $J$ ($\bar s$ does not depend on $\mu$: $\bar\tau=\mathcal I_{2,0}$ has $\beta=0$). Now let $J\to\infty$ at fixed $k$. By \cref{lem:s10:V}, $\alpha^{\mathrm{tor}}_{k,T_J}(M)\to\alpha^\infty_k(M)$ for $M\in\{0.7,1\}$, and its proof (\cref{ssec:a06:V}) establishes the convergence of the recursion \eqref{eq:s14:CTdef} order by order: $\dm_{(\alpha\beta),T_J}(\eps_k;M)\to\dm_{(\alpha\beta)}(\eps_k;M)$ for each of the five orders (alternatively: a polynomial of bounded degree in $(e^2,\lambda)$ that converges at every point of an open set converges coefficient by coefficient, and \cref{lem:s10:V} holds at every point of King's range). So $S^{T_J}_{\alpha\beta,k}\to S_{\alpha\beta,k}$, and closed inequalities pass to the limit.
\end{proof}

\begin{lemma}[the torus floor can be removed]\label{lem:s14:TF}
For every $\lmu\in(0,1]$, $n\ge1$ and every $k$ with $\eps_k\le1$, the four inequalities of \cref{lem:s14:E} hold with $\bar X_n(\lmu,\infty)$, the upper sums formed with the $\eps\Z^3$ majorants $\mathcal I_{\alpha,\beta}(M,\lmu;\infty)=c_NC_3(M,\lmu)$ and with $\bar\omega=w$.
\end{lemma}

\begin{proof}
Fix $k$, $\lmu$, $n$. By \cref{lem:s14:E}, $2S_{\alpha\beta,k}\ge-\bar X_n(\lmu,\ell_0)$ for \emph{every} $\ell_0>0$: the torus floor is a free parameter of (G4), and the tori $T_J$ of the volume limit eventually exceed any $\ell_0$. It therefore suffices that $\bar X_n(\lmu,\ell_0)\to\bar X_n(\lmu,\infty)$ as $\ell_0\to\infty$. Each $\bar X_n$ is a finite sum of cell values, and each cell value is a fixed finite expression (sums, products and minima, with the fixed prefactors) in the majorants $\mathcal I_{\alpha,\beta}(a,\lmu;\ell_0)$ at the left endpoint $a$ of the cell and in $\bar\omega(\lmu,\ell_0)$, hence continuous in these. As $\ell_0\to\infty$, $\mathcal I_{\alpha,\beta}(a,\lmu;\ell_0)\to c_NC_3(a,\lmu)$ by \eqref{eq:s14:Iab}, since $C_2$ and $C_1$ are finite numbers that do not depend on $\ell_0$; and $0\le\hat E(\lmu,\ell_0)\le\bar E(\lmu,\ell_0)\to0$, each of the four terms of $\bar E$ being a positive constant times $\ell_0^{-1}$, $\ell_0^{-1}E_1(\lmu^2\ell_0^2/8)$, $\ell_0^{-2}\mathrm{erfc}(\lmu\ell_0/(2\sqrt2))$ or $\ell_0^{-3}e^{-\lmu^2\ell_0^2/8}$, so $\bar\omega\to w$.
\end{proof}

\begin{lemma}[monotonicity in the photon mass; interval bounds]\label{lem:s14:MON}
Let $0<\mu_{\rm lo}\le\mu_{\rm hi}\le1$. Define the \emph{interval upper sum} $\bar X^{\rm int}_n([\mu_{\rm lo},\mu_{\rm hi}])$ as $\bar X_n(\cdot,\infty)$ with every majorant evaluated at $\mu=\mu_{\rm lo}$ and, on the cell $[a,b]$, the prefactor $|\mu^2-2M^2|$ replaced by $\max\{|\mu^2-2M^2|:\mu\in\{\mu_{\rm lo},\mu_{\rm hi}\},\,M\in\{a,b\}\}$. Then $\bar X_n(\mu,\infty)\le\bar X^{\rm int}_n([\mu_{\rm lo},\mu_{\rm hi}])$ for every $\mu\in[\mu_{\rm lo},\mu_{\rm hi}]$, so the inequalities of \cref{lem:s14:TF} hold with $\bar X^{\rm int}_n([\mu_{\rm lo},\mu_{\rm hi}])$ for all $\lmu\in[\mu_{\rm lo},\mu_{\rm hi}]$ at once. Moreover, if $\mu_{\rm hi}\le0.989949$, the interval upper sum equals the point upper sum $\bar X_n(\mu_{\rm lo},\infty)$.
\end{lemma}

\begin{proof}
A cell value is non-decreasing in each majorant and in the prefactor (sums, products and minima of non-negative quantities with non-negative coefficients). (i) Each admissible closed form of $C_3$ is non-increasing in $b=\mu$: the $\nu$-form, since $\nu=\min(a,b)$ is non-decreasing in $b$ and $3-2g<0$; the Beta form in $a$ is $b$-free; the Beta form in $b$ carries $b^{2c'-2\beta}$ with $c'=\frac{3-2\alpha}2<\beta$, as $g>\frac32$; the H\"older form carries $b^{(3-2g)\beta/g}$; the exact forms $\frac1{4\pi(a+b)}$ and $\frac1{8\pi a(a+b)^2}$ decrease in $b$. A minimum of non-increasing functions is non-increasing, so every $\mathcal I_{\alpha,\beta}(a,\mu;\infty)$ is non-increasing in $\mu$, and $\bar\omega=w$ does not depend on $\mu$. (ii) $\mu\mapsto|\mu^2-2M^2|$ is convex in $\mu^2$ and $M\mapsto|\mu^2-2M^2|$ convex in $M^2$, so the maximum over the rectangle $[\mu_{\rm lo},\mu_{\rm hi}]\times[a,b]$ is attained at a corner. Hence every ingredient of the cell value at $\mu$ is at most the corresponding ingredient of the interval cell value. For the last claim: on $[0.7,1]$ one has $2M^2\ge0.98>0.989949^2$, so $|\mu^2-2M^2|=2M^2-\mu^2$ on the whole rectangle, maximal at $(\mu_{\rm lo},b)$, which is also the point prefactor at $\mu_{\rm lo}$. (For the interval $[0.98,0.99]$ the corner $(0.99,0.7)$ of the first cell gives $|0.9801-0.98|=10^{-4}$, below the value $2\cdot0.70075^2-0.98^2=0.0217$ at $(0.98,0.70075)$, so the equality holds there too; only the interval $[0.99,1]$ differs, by $8.5786210$ against the point value $8.5785880$ of $\bar B_3$.)
\end{proof}

\begin{lemma}[a closed form, uniform in the photon mass]\label{lem:s14:CF}
For $M\in[0.7,1]$ put $\bar\tau^*(M)\defeq\mathcal I_{2,0}(M;\infty)=\frac{c_N}{8\pi M}$, $\bar\psi^*(M)\defeq\frac{c_N}{4\pi M}$, and for $X\in\{\Theta,J_s,J'_s\}$ let $X^*(M)$ be $c_N$ times the minimum over the $\mu$-free admissible forms of $C_3$ only (the Beta forms in $a=M$ and, for $\beta=0$, the H\"older form $\gamma_3(g)M^{3-2g}$; the $\nu$-forms, the Beta forms in $b$, the exact form at $(2,1)$ and the H\"older forms with $\beta>0$ are dropped, and the exact form at $(1,1)$ is replaced by its bound $\frac1{4\pi a}$, which is $\bar\psi^*$). Put $\bar\beta^*\defeq2M[2\bar\tau^*+8\bar\psi^*]$,
\begin{gather*}
\varphi_3^*(M)\defeq4M\big[2\sqrt2\,\bar\Theta^*\bar J^*_{3/2}+2M^2\bar\psi^*\bar J^*_0+3\sqrt3\,w\,\bar J^*_1\big]+\bar\beta^*\bar\tau^*,\qquad
\varphi_A^*(M)\defeq6M\big[2\sqrt2\,\bar\Theta^*\bar J'^*_{3/2}+2M^2\bar\psi^*\bar J'^*_0+3\sqrt3\,w\,\bar J'^*_1\big]+\bar\beta^*\bar\psi^*,\\
P(M)\defeq\varphi_A^*+52M\bar\psi^{*2}+16M\bar\psi^*\bar\tau^*,\qquad Q(M)\defeq\frac{c_N}{8\pi}M\big(28\bar\tau^*+40\bar\psi^*\big).
\end{gather*}
Then $16\varphi_3(M;\mu,\infty)\le16\varphi_3^*(M)$ and $b_4(M;\mu,\infty)\le P(M)+Q(M)/\mu$ for all $M\in[0.7,1]$ and $\mu\in(0,1]$. Consequently, with the $400$-cell upper sums $B_3^*\defeq\sum h\,(16\varphi_3^*)^{\rm cell}$, $P^*\defeq\sum hP^{\rm cell}$, $Q^*\defeq\sum hQ^{\rm cell}$ (majorants at the left endpoint, the prefactors $M$ and $2M^2$ at the right endpoint), for every $\lmu\in(0,1]$ and every $k$ with $\eps_k\le1$:
\[
2S_{21,k}\ge-B_3^*,\qquad 2S_{40,k}\ge-\big(P^*+Q^*/\lmu\big),\qquad B_3^*=13.187549,\quad P^*=4.7629000,\quad Q^*=0.77120851 .
\]
\end{lemma}

\begin{proof}
Each starred majorant is an upper bound of the corresponding majorant at $(M,\mu;\infty)$ for every $\mu>0$: a minimum over a subset of the admissible forms is at least the minimum over all of them, and $\mathcal I_{1,1}(M,\mu;\infty)\le c_N\frac1{4\pi(M+\mu)}\le\frac{c_N}{4\pi M}$. Further $\bar\delta(\mu)=\mathcal I_{0,2}(M,\mu;\infty)\le c_N\frac1{4\pi^2}B(\frac32,\frac12)\mu^{-1}=\frac{c_N}{8\pi\mu}$ (the Beta form in $b$, admissible since $\alpha=0<\frac32$ and $\beta=2>\frac32$); $|\mu^2-2M^2|\le2M^2$, since $\mu^2\le1\le4M^2$; $\bar\omega=w$ (\cref{lem:s14:TF}); and in $\varphi_{D2}$ the minimum is at most its first branch. Substituting these bounds into the formulas (G3), which are non-decreasing in every majorant, gives $16\varphi_3\le16\varphi_3^*$ and
\[
b_4\le\varphi_A^*+28M\bar\tau^*\frac{c_N}{8\pi\mu}+2M\bar\psi^*\Big(2\bar\psi^*+\frac{4c_N}{8\pi\mu}\Big)+8M\bar\psi^*\Big(6\bar\psi^*+2\bar\tau^*+\frac{4c_N}{8\pi\mu}\Big)=P(M)+\frac{Q(M)}\mu .
\]
The starred majorants are non-increasing in $M$ (powers $M^{2c-2\alpha}$ with $c<\alpha$, and $1/M$), so the cell rule \eqref{eq:s14:cell} applies to $16\varphi_3^*$, $P$ and $Q$ with the prefactors $M$, $2M^2$ at the right endpoint. The two inequalities follow \emph{cellwise} from \cref{lem:s14:TF}: for a cell $[a,b]$ the cell value of $16\varphi_3$ (resp.\ $b_4$) at $(\lmu,\infty)$ is a non-decreasing expression in the majorants at $(a,\lmu;\infty)$, in $\bar\omega$ and in the prefactor, each of which is at most its starred counterpart at $a$ (with $\bar\delta\le c_N/(8\pi\lmu)$, $\max_{M\in\{a,b\}}|\lmu^2-2M^2|\le2b^2$, $\bar\omega=w$); so the cell value is at most the cell value of $16\varphi_3^*$ (resp.\ of $P+Q/\lmu$). Summing over the $400$ cells, $\bar B_{3,400}(\lmu,\infty)\le B_3^*$ and $\bar B^-_{4,400}(\lmu,\infty)\le P^*+Q^*/\lmu$, and \cref{lem:s14:TF} gives the two inequalities. The three numbers are the $400$-cell sums, computed as described after \cref{tab:s14:T3}. (The slope bounds (G3) are not invoked with the starred majorants in place of the one-loop majorants: as printed they make only $\bar\tau$, $\bar\psi$ generic; the domination above is all that is used.)
\end{proof}

% s14_t3_table: the interval upper sums at ell_0 = infinity and the derived region numbers (input from sec/s14.tex)

\subsection*{The numbers}

Let $\hat\mu=\lmu/m\in(0,1]$. The grid is $\{1,0.99,0.98,0.95,0.9,0.85,\dots,0.15,0.1,0.08,0.06,0.05,0.04,0.03,0.02,0.015,0.01,0.005,0.002,0.001,0.0005,0.0002,0.0001\}$ ($35$ points), and the \emph{rows} of \cref{tab:s14:T3} are the $34$ intervals $[\mu_{\rm lo},\mu_{\rm hi}]$ between consecutive grid points. For $\hat\mu\in[10^{-4},1]$ the \emph{row of $\hat\mu$} is the row containing $\hat\mu$ (at a grid point, where two rows contain $\hat\mu$, the row with $\mu_{\rm lo}=\hat\mu$), and $\bar B_3(\hat\mu)$, $\bar B_4^-(\hat\mu)$, $\bar j(\hat\mu)$ denote the interval upper sums $\bar B^{\rm int}_{3,400}$, $\bar B^{-,{\rm int}}_{4,400}$, $\bar j^{\rm int}_{400}$ of that row (\cref{lem:s14:MON}), and $\bar s\defeq\bar s_{400}=6.5231964$; for $\hat\mu<10^{-4}$, below the table, we read $\bar B_3(\hat\mu)\defeq B_3^*$, $\bar B_4^-(\hat\mu)\defeq P^*+Q^*/\hat\mu$ and $\bar j(\hat\mu)\defeq0.092610085$ (\cref{lem:s14:CF} for the first two; the last is the $\mu$-free branch $\int2M\bar\tau$ of $\bar j$, see the proof of \cref{prop:s14:T3}(i)). With $A\defeq\frac12(1-M_1^2)=0.255$, $b\defeq\frac{4(1-M_1)}{2\pi}=0.1909859$, $L\defeq\frac{4\ln(1/M_1)}{\pi^2}=0.1445549$ (\cref{lem:s14:A,lem:s14:B,lem:s14:C} at $N=2$, $M_1=0.7$), $I_\infty(\hat\mu)=\frac{0.7+1.7\hat\mu}{4\pi\cdot1.7\,(\hat\mu+0.7)(1+\hat\mu)}$ (so $I_\infty(1)=0.0330425$) and $I^*=\frac1{4\pi\cdot1.7}=0.0468103$, put
\begin{equation}\label{eq:s14:sigmalow}
\sigma_{\rm low}(\hat\lambda,\hat e^2;\hat\mu)\defeq A\big(1-\hat e^2I_\infty(\mu_{\rm lo})\big)+b\hat\lambda-L\hat\lambda^2-\tfrac12\bar B_3(\hat\mu)\,\hat e^2\hat\lambda-\tfrac12\bar B_4^-(\hat\mu)\,\hat e^4\qquad(\hat\mu\in[10^{-4},1]),
\end{equation}
\begin{equation}\label{eq:s14:sigmastar}
\sigma^*(\hat\lambda,\hat e^2;\hat\mu)\defeq A\big(1-\hat e^2I^*\big)+b\hat\lambda-L\hat\lambda^2-\tfrac12B_3^*\,\hat e^2\hat\lambda-\tfrac12\big(P^*+Q^*/\hat\mu\big)\hat e^4\qquad(\hat\mu\in(0,1]),
\end{equation}
and $\sigma_{\rm low}(\cdot;\hat\mu)\defeq\sigma^*(\cdot;\hat\mu)$ for $\hat\mu<10^{-4}$. The derived columns of the table are: $e^2_{\max}$, the positive root of $\sigma_{\rm low}(0,\cdot;\hat\mu)$, $e^2_{\max}=\big(-AI_\infty+\sqrt{A^2I_\infty^2+2A\bar B_4^-}\big)/\bar B_4^-$; the side $c$ of the largest box $[0,c]^2$ on which $\sigma_{\rm low}\ge0$; and the value $\sigma_{\rm low}(0.1,0.1;\hat\mu)$, where a negative entry means only that the bound does not place $[0,0.1]^2$ in the region.

\begin{table}[htbp]
\centering\footnotesize
\setlength{\tabcolsep}{4pt}
\begin{tabular}{l|lll|l|lll}
$[\mu_{\rm lo},\mu_{\rm hi}]$ & $\bar j$ & $\bar B_3$ & $\bar B_4^-$ & $AI_\infty(\mu_{\rm lo})$ & $e^2_{\max}(\lambda=0)$ & $c$ & $\sigma_{\rm low}(0.1,0.1)$\\\hline
$[0.99,1]$ & 0.085313936 & 8.5786210 & 1.5854190 & 0.008458 & 0.5618 & 0.2390 & 0.2209\\
$[0.98,0.99]$ & 0.085757165 & 8.6036471 & 1.6023646 & 0.008491 & 0.5588 & 0.2385 & 0.2207\\
$[0.95,0.98]$ & 0.087020288 & 8.6804035 & 1.6554211 & 0.008589 & 0.5498 & 0.2369 & 0.2201\\
$[0.9,0.95]$ & 0.088887935 & 8.8139269 & 1.7519811 & 0.008757 & 0.5345 & 0.2341 & 0.2189\\
$[0.85,0.9]$ & 0.090427040 & 8.9551127 & 1.8603008 & 0.008930 & 0.5188 & 0.2313 & 0.2176\\
$[0.8,0.85]$ & 0.091595352 & 9.1048479 & 1.9826438 & 0.009108 & 0.5026 & 0.2283 & 0.2163\\
$[0.75,0.8]$ & 0.092342924 & 9.2641791 & 2.1218953 & 0.009291 & 0.4858 & 0.2251 & 0.2147\\
$[0.7,0.75]$ & 0.092610085 & 9.4343534 & 2.2812536 & 0.009480 & 0.4686 & 0.2217 & 0.2131\\
$[0.65,0.7]$ & 0.092610085 & 9.6168755 & 2.4625429 & 0.009673 & 0.4511 & 0.2182 & 0.2112\\
$[0.6,0.65]$ & 0.092610085 & 9.8135862 & 2.6698136 & 0.009871 & 0.4333 & 0.2145 & 0.2092\\
$[0.55,0.6]$ & 0.092610085 & 10.026776 & 2.9084862 & 0.010073 & 0.4152 & 0.2105 & 0.2069\\
$[0.5,0.55]$ & 0.092610085 & 10.259350 & 3.1856310 & 0.010279 & 0.3969 & 0.2063 & 0.2044\\
$[0.45,0.5]$ & 0.092610085 & 10.515087 & 3.5155645 & 0.010488 & 0.3779 & 0.2017 & 0.2014\\
$[0.4,0.45]$ & 0.092610085 & 10.799045 & 3.9183359 & 0.010697 & 0.3580 & 0.1967 & 0.1979\\
$[0.35,0.4]$ & 0.092610085 & 11.118258 & 4.4224644 & 0.010906 & 0.3371 & 0.1912 & 0.1938\\
$[0.3,0.35]$ & 0.092610085 & 11.477171 & 5.0733664 & 0.011111 & 0.3148 & 0.1850 & 0.1887\\
$[0.25,0.3]$ & 0.092610085 & 11.851059 & 5.9437544 & 0.011309 & 0.2910 & 0.1781 & 0.1825\\
$[0.2,0.25]$ & 0.092610085 & 12.193509 & 7.0632067 & 0.011495 & 0.2670 & 0.1709 & 0.1752\\
$[0.15,0.2]$ & 0.092610085 & 12.433452 & 8.5440078 & 0.011662 & 0.2429 & 0.1635 & 0.1665\\
$[0.1,0.15]$ & 0.092610085 & 12.659419 & 11.303445 & 0.011801 & 0.2113 & 0.1525 & 0.1516\\
$[0.08,0.1]$ & 0.092610085 & 12.756284 & 13.316343 & 0.011846 & 0.1948 & 0.1460 & 0.1411\\
$[0.06,0.08]$ & 0.092610085 & 12.857230 & 16.620716 & 0.011884 & 0.1744 & 0.1370 & 0.1240\\
$[0.05,0.06]$ & 0.092610085 & 12.909321 & 19.239396 & 0.011899 & 0.1621 & 0.1310 & 0.1107\\
$[0.04,0.05]$ & 0.092610085 & 12.962542 & 23.145186 & 0.011912 & 0.1479 & 0.1233 & 0.0909\\
$[0.03,0.04]$ & 0.092610085 & 13.016932 & 29.623519 & 0.011923 & 0.1308 & 0.1132 & 0.0582\\
$[0.02,0.03]$ & 0.092610085 & 13.072534 & 42.530531 & 0.011931 & 0.1092 & 0.0987 & $-0.0066$\\
$[0.015,0.02]$ & 0.092610085 & 13.100802 & 55.411533 & 0.011933 & 0.0957 & 0.0887 & $-0.0712$\\
$[0.01,0.015]$ & 0.092610085 & 13.129390 & 81.146532 & 0.011935 & 0.0791 & 0.0753 & $-0.2000$\\
$[0.005,0.01]$ & 0.092610085 & 13.158304 & 158.29597 & 0.011937 & 0.0566 & 0.0555 & $-0.5859$\\
$[0.002,0.005]$ & 0.092610085 & 13.175811 & 389.67594 & 0.011937 & 0.0361 & 0.0360 & $-1.7428$\\
$[0.001,0.002]$ & 0.092610085 & 13.181673 & 775.28604 & 0.011937 & 0.0256 & 0.0256 & $-3.6709$\\
$[0.0005,0.001]$ & 0.092610085 & 13.184609 & 1546.4975 & 0.011937 & 0.0181 & 0.0181 & $-7.5270$\\
$[0.0002,0.0005]$ & 0.092610085 & 13.186373 & 3860.1248 & 0.011937 & 0.0114 & 0.0114 & $-19.095$\\
$[0.0001,0.0002]$ & 0.092610085 & 13.186961 & 7716.1679 & 0.011937 & 0.0081 & 0.0081 & $-38.375$
\end{tabular}
\caption{The interval upper sums $\bar j$, $\bar B_3$, $\bar B_4^-$ of \cref{lem:s14:MON} on the $400$-cell partition of $[0.7,1]$ at $\ell_0=\infty$ (\cref{lem:s14:TF}), valid for every $\hat\mu$ in the row; $\bar s=6.5231964$ on every row. The three upper sums are rounded upward to eight significant digits; the derived columns $e^2_{\max}$, $c$ are rounded down and $\sigma_{\rm low}(0.1,0.1)$ is rounded down. Units $m=1$, $N=2$, $M_1=0.7$.}\label{tab:s14:T3}
\end{table}

\paragraph{How the entries were computed, and what is exact.} The upper sums are the finite sums \eqref{eq:s14:cell} over the $400$ cells of the formulas (G2), (G3), with the interval rule of \cref{lem:s14:MON}; their validity as upper bounds is the inequality \eqref{eq:s14:cell}, which is exact once the finitely many closed forms are evaluated correctly. At $\ell_0=\infty$ only the $c_NC_3$ term of \eqref{eq:s14:Iab} and $\bar\omega=w\le0.25273101$ enter, so no series and no special function beyond $\Gamma$ and $B$ is involved. The evaluation was carried out in floating-point arithmetic with $50$ significant digits; every final upper sum was multiplied by $1+10^{-30}$ and rounded upward to eight significant digits, the derived region numbers were computed from the rounded upper sums with coefficients rounded up and sizes rounded down, and a second run at $30$ digits agrees with the first to all printed digits. The rounding of the elementary and Gamma functions at $50$ digits is many orders below the safety factor and the final rounding; interval validity rests on \cref{lem:s14:MON} and finitely many evaluations, not on interval arithmetic. The same evaluation at $\hat\mu=1$ and $\ell_0=96$, with $\bar\omega=w+\min\{\bar E,\tilde E\}$ as in (G2), reproduces the four numbers $0.0886293$, $7.179064$, $9.330427$, $1.703784$ of \cite[Cor.~G.4]{SymPhase} to every printed digit (it gives $0.0886292849$, $7.17906354$, $9.3304261$, $1.70378338$); this checks the transcription of the formulas and of the set of admissible forms against the torus paper's printed certificate, which exercises every form that is active at $\mu=1$. Internal consistency of the table: $\bar s_{400}=6.5231964$ at every $\hat\mu$ ($\bar s$ is $\mu$-free); $\bar j=0.092610085$ for all $\hat\mu\le0.7$, the $\mu$-free branch $\int2M\bar\tau$ of $\bar j$ (analytically $\frac{c_N}{4\pi}(1-M_1)=0.09253$ before the cell rounding); the interval values equal the point values at $\mu_{\rm lo}$ except on $[0.99,1]$ (\cref{lem:s14:MON}); $\bar B^-_{4,400}(\hat\mu,\infty)\le P^*+Q^*/\hat\mu$ at every grid point ($7716.17\le7716.85$ at $\hat\mu=10^{-4}$, $1.5854\le5.534$ at $\hat\mu=1$); and $\hat\mu\,\bar B_4^-(\hat\mu)\to0.7716$ at $\hat\mu=10^{-4}$ against $Q^*=0.7712$. The closed form of \cref{lem:s14:CF} is weaker than the table near $\hat\mu=1$ (it uses $\bar\psi^*$, $|\mu^2-2M^2|\le2M^2$ and the first branch of the minimum) and is within one per cent of it for $\hat\mu\le10^{-2}$.


\begin{proposition}[the explicit coupling region of \Tthree\ at every photon mass]\label{prop:s14:T3}
Let $N=2$, $M_1=0.7m$ and $\lmu\in(0,m]$; write $\hat\lambda=\lambda/m$, $\hat e^2=e^2/m$, $\hat\mu=\lmu/m$. Let $\dm_k(M)$ be King's counterterm \eqref{eq:s14:CTdef} on $\Lk$ at the photon mass $\lmu$, $S_{\alpha\beta,k}$ the coefficients of \cref{lem:s14:struct}, $\sigma_1(M_1)=\liminf_ks^\infty_k(M_1)$ as in \cref{hyp:s03:T2}, and let $\sigma_{\rm low}$, $\sigma^*$, $\bar j(\hat\mu)$, $\bar s$, $\bar B_3(\hat\mu)$, $\bar B_4^-(\hat\mu)$, $A$, $b$, $L$, $I_\infty$, $I^*$ be as defined before \cref{tab:s14:T3}, $B_3^*$, $P^*$, $Q^*$ as in \cref{lem:s14:CF}, and $\lambda_c\defeq\frac{b+\sqrt{b^2+4AL}}{2L}=2.14398\ldots$.
\begin{enumerate}[label=\textup{(\roman*)},leftmargin=*]
\item \emph{(The five orders.)} For every $k$ with $\eps_k\le1/m$, in units $m=1$: $S_{01,k}>0$; $2S_{20,k}\ge-\bar j(\hat\mu)\ge-0.092611$; $2S_{02,k}\ge-\bar s=-6.5231964$; $2S_{21,k}\ge-\bar B_3(\hat\mu)$ and $2S_{21,k}\ge-B_3^*$; $2S_{40,k}\ge-\bar B_4^-(\hat\mu)$ and $2S_{40,k}\ge-(P^*+Q^*/\hat\mu)$.
\item \emph{(The lower bound.)} $\sigma_1(0.7m)\ge m^2\max\big\{\sigma_{\rm low}(\hat\lambda,\hat e^2;\hat\mu),\ \sigma^*(\hat\lambda,\hat e^2;\hat\mu)\big\}$.
\item \emph{(The region.)} \Tthree\ holds on $R(\lmu)\defeq\{(\lambda,e^2):\sigma_{\rm low}(\hat\lambda,\hat e^2;\hat\mu)>0\}$, with $\sigma_1(0.7m)\ge m^2\sigma_{\rm low}>0$ there. $R(\lmu)$ contains the box $[0,c(\hat\mu)m]^2$ and the segment $\{\lambda=0,\ e^2<e^2_{\max}(\hat\mu)m\}$, with $c$ and $e^2_{\max}$ from \cref{tab:s14:T3} for $\hat\mu\in[10^{-4},1]$, and the segment $\{e=0,\ \lambda<\lambda_cm\}$ for every $\hat\mu$. The map $\hat\mu\mapsto\sigma_{\rm low}(\hat\lambda,\hat e^2;\hat\mu)$ is non-decreasing (on $[10^{-4},1]$ a step function of $\hat\mu$, constant on each row of the table; below $10^{-4}$ the closed form $\sigma^*$), so $R(\lmu)\supseteq R(\lmu')$ for $\lmu\ge\lmu'$. In particular $[0,0.2m]^2\subseteq R(\lmu)$ for $\lmu\in[0.5m,m]$; $[0,0.15m]^2\subseteq R(\lmu)$ for $\lmu\in[0.1m,m]$; $[0,0.1m]^2\subseteq R(\lmu)$ for $\lmu\in[0.03m,m]$, with $\sigma_1(0.7m)\ge0.0582\,m^2$ at $(\lambda,e^2)=(0.1m,0.1m)$; $[0,0.075m]^2\subseteq R(\lmu)$ for $\lmu\in[0.01m,m]$; and at $\lmu=m$, $R(m)$ contains the box $[0,0.239m]^2$ and the segment $\{\lambda=0,\ e^2<0.5618\,m\}$, and $\sigma_1(0.7m)\ge0.2209\,m^2$ on $[0,0.1m]^2$.
\item \emph{(Uniform in the photon mass; closed form.)} \Tthree\ holds whenever $\sigma^*(\hat\lambda,\hat e^2;\hat\mu)>0$, with $\sigma_1(0.7m)\ge m^2\sigma^*$; in particular, for every $\lmu\in(0,m]$, on each of the boxes $\{\lambda\le\lambda_1m,\ e^2\le c_1\sqrt{\lmu m}\}$ with $(\lambda_1,c_1)\in\{(0,0.3014),(0.05,0.2533),(0.1,0.2152),(0.2,0.1616),(0.5,0.0884),(1,0.0447)\}$.
\item \emph{(Every $k$.)} For every $k$ with $\eps_k\le1/m$: $s^\infty_k(0.7m)\ge m^2\big[A-\frac12\bar j(\hat\mu)\hat e^2-\frac12\bar s\hat\lambda^2-\frac12\bar B_3(\hat\mu)\hat e^2\hat\lambda-\frac12\bar B_4^-(\hat\mu)\hat e^4\big]$.
\item \emph{($e=0$; exact.)} At $e=0$, for every $M_1\in(0,m)$, every $\lmu>0$ and every $N\ge1$: $\lim_ks^\infty_k(M_1)=\frac12(m^2-M_1^2)+\frac{(N+2)\lambda(m-M_1)}{2\pi}-\frac{(N+2)\lambda^2}{\pi^2}\ln\frac m{M_1}$ exactly, and \Tthree\ holds if and only if $\lambda<\frac\pi4\big(1+\sqrt{1+\frac{16}{N+2}}\big)m$, which is $\frac\pi4(1+\sqrt5)\,m=2.5416\ldots\,m$ for $N=2$; at $M_1=0.7m$, $N=2$, the exact window is $\lambda<\lambda_cm$, $\lambda_c=2.14398\ldots$, in particular $\lambda\le2.1439\,m$.
\item \emph{(Two-sided.)} On each region of (iii), (iv) and (vi) the two-sided clause \sxivNE\ holds as well, with no use of \cref{prop:s10:Upm}.
\end{enumerate}
\end{proposition}

Outside these regions the sign clause \Tthree\ is not decided here. The region is explicit at every photon mass but not uniform in it: by \eqref{eq:s14:sigmastar}, $e^2_{\max}(\hat\mu)\sim\sqrt{2A/Q^*}\sqrt{\hat\mu}=0.81\sqrt{\hat\mu}$ as $\hat\mu\to0$ (\cref{tab:s14:T3}: $\hat\mu\,\bar B_4^-(\hat\mu)\to0.77$), and $\bigcap_{\lmu\in(0,m]}R(\lmu)=\{e=0,\ \lambda<\lambda_cm\}$, since for $e>0$ the term $-\frac12(Q^*/\hat\mu)\hat e^4$ of $\sigma^*$ dominates as $\hat\mu\to0$. The $1/\hat\mu$ sits in the order-$e^4$ majorant alone: the low orders and $B_3^*$ are uniform in the photon mass.

% s14_t3_proof: the proof of prop:s14:T3 (input from sec/s14.tex)
\begin{proof}[Proof of \cref{prop:s14:T3}]
Units $m=1$, $N=2$, $M_1=0.7$; $\hat\mu=\lmu\in(0,1]$.

\emph{(i).} $S_{01,k}>0$ is \cref{lem:s14:A}(i). The four bounds with the interval upper sums are \cref{lem:s14:TF} with \cref{lem:s14:MON}, applied to the row of $\hat\mu$ (for $\hat\mu\ge10^{-4}$); the bounds with $B_3^*$ and $P^*+Q^*/\hat\mu$ are \cref{lem:s14:CF}; for $\hat\mu<10^{-4}$ the first bound is the $\mu$-free branch of $\bar j$: $j\le2M\bar\tau$ gives $\bar j_{400}(\hat\mu,\infty)\le\sum h\,(2M\bar\tau^*)^{\rm cell}=0.092610085$ for every $\hat\mu$, by the cell rule \eqref{eq:s14:cell}. The numbers $\bar j\le0.092610085<0.092611$ and $\bar s=6.5231964$ are read off \cref{tab:s14:T3}. The interval value $\bar B_4^-(\hat\mu)$, taken at $\mu_{\rm lo}$, and the closed form $P^*+Q^*/\hat\mu$, taken at $\hat\mu$, are two separate bounds, neither of which dominates the other on a whole row.

\emph{(ii).} By \eqref{eq:s14:struct} and \cref{lem:s14:A}(iii), \cref{lem:s14:B}(iii), \cref{lem:s14:C}(iii), the first four coefficients converge, so
\[
\sigma_1(0.7)=\liminf_ks^\infty_k(0.7)=A+b\hat\lambda-\hat e^2AI_\infty(\hat\mu)-L\hat\lambda^2+\liminf_k\big(\hat e^2\hat\lambda S_{21,k}+\hat e^4S_{40,k}\big),
\]
and by (i), for all but finitely many $k$ (those with $\eps_k\le1$), $\hat e^2\hat\lambda S_{21,k}+\hat e^4S_{40,k}\ge-\frac12\bar B_3(\hat\mu)\hat e^2\hat\lambda-\frac12\bar B_4^-(\hat\mu)\hat e^4$; with $I_\infty(\hat\mu)\le I_\infty(\mu_{\rm lo})$ (\cref{lem:s14:B}(iv)) this is $\sigma_1(0.7)\ge\sigma_{\rm low}$. The same argument with the closed-form bounds of (i) and $I_\infty(\hat\mu)\le I^*$ gives $\sigma_1(0.7)\ge\sigma^*$. (Cellwise, the interval value of $16\varphi_3$ at $\mu_{\rm lo}$ is at most the starred cell value, so $\bar B_3(\hat\mu)\le B_3^*$ always holds; for $b_4$ the interval value carries $\bar\delta(\mu_{\rm lo})\approx c_N/(8\pi\mu_{\rm lo})$ while the closed form carries $Q^*/\hat\mu$, so neither bound dominates the other on a whole row; hence the maximum in (ii).)

\emph{(iii).} The first sentence is (ii) and the definition of \Tthree\ (\cref{hyp:s03:T3}). $\sigma_{\rm low}$ is concave in $\hat\lambda$ at fixed $\hat e^2$ (the term $-L\hat\lambda^2$) and concave in $\hat e^2$ at fixed $\hat\lambda$ (the term $-\frac12\bar B_4^-\hat e^4$), so its minimum over a box $[0,c]^2$ is attained at a corner: the minimum over $\hat\lambda\in[0,c]$ of a concave function is attained at an endpoint, the two resulting functions of $\hat e^2$ are concave and so is their minimum, whose minimum over $[0,c]$ is again at an endpoint. The side $c(\hat\mu)$ of \cref{tab:s14:T3} is the largest $c$ for which the four corner values are $\ge0$, rounded down; $e^2_{\max}$ is the positive root of $\sigma_{\rm low}(0,\cdot)=A(1-\hat e^2I_\infty)-\frac12\bar B_4^-\hat e^4$, rounded down; and at $e=0$, $\sigma_{\rm low}=A+b\hat\lambda-L\hat\lambda^2>0$ if and only if $\hat\lambda<\lambda_c=\frac{b+\sqrt{b^2+4AL}}{2L}=2.14398\ldots$ (as in (vi)). \emph{Monotonicity in $\hat\mu$:} on $[10^{-4},1]$, $\sigma_{\rm low}(\cdot;\hat\mu)$ is a step function of $\hat\mu$, built row by row from $I_\infty(\mu_{\rm lo})$ and the interval values $\bar B_3$, $\bar B_4^-$ of the row, which enter with negative coefficients; so it suffices that these three row values do not decrease from each row to the row below it. For every row but the top one the interval values are the point values at $\mu_{\rm lo}$ (\cref{lem:s14:MON}, last clause, and the remark at the end of its proof for the row $[0.98,0.99]$), the point values are non-increasing in $\mu$ (\cref{lem:s14:MON}(i), with the point prefactor $\max_{M\in\{a,b\}}|\mu^2-2M^2|=2b^2-\mu^2$ non-increasing in $\mu$ for $\mu\le0.989949$) and so is $I_\infty$ (\cref{lem:s14:B}(iv)); hence the values of a row, taken at its $\mu_{\rm lo}$, are at most those of the row below, whose left end is smaller. For the top row $[0.99,1]$ this is read off the table ($8.5786210\le8.6036471$, $1.5854190\le1.6023646$, $0.008458\le0.008491$); by inspection each of the three columns is non-decreasing down the table through all $34$ rows. Across $\hat\mu=10^{-4}$: $\sigma^*$ is non-decreasing in $\hat\mu$, and $\sigma^*(\cdot;\hat\mu)\le\sigma_{\rm low}(\cdot;10^{-4})$ for $\hat\mu<10^{-4}$ because $I^*\ge I_\infty(10^{-4})$, $B_3^*\ge\bar B_3(10^{-4})=13.186961$ and $P^*+Q^*/\hat\mu>7716.84>\bar B_4^-(10^{-4})=7716.1679$. At a grid point both rows containing it give valid bounds, and the convention selects the row with $\mu_{\rm lo}=\hat\mu$. Hence $\hat\mu\mapsto\sigma_{\rm low}(\hat\lambda,\hat e^2;\hat\mu)$ is non-decreasing and $R(\lmu)\supseteq R(\lmu')$ for $\lmu\ge\lmu'$. \emph{The listed boxes:} $c\ge0.2063$ on the row $[0.5,0.55]$, $c\ge0.1525$ on $[0.1,0.15]$, $c\ge0.1132$ and $\sigma_{\rm low}(0.1,0.1)\ge0.0582$ on $[0.03,0.04]$, $c\ge0.0753$ on $[0.01,0.015]$ (\cref{tab:s14:T3}), together with the monotonicity; at $\hat\mu=1$ the row $[0.99,1]$ gives $c\ge0.2390$ and $e^2_{\max}\ge0.5618$.

\emph{(iv).} $\sigma^*$ is concave in each variable separately, so again the corners decide. Let $(\hat\lambda,\hat e^2)$ lie in the box $\{\hat\lambda\le\lambda_1,\ \hat e^2\le c_1\sqrt{\hat\mu}\}$ with $\hat\mu\le1$. At a corner $(\hat\lambda,c_1\sqrt{\hat\mu})$, $\hat\lambda\in\{0,\lambda_1\}$,
\[
\sigma^*=A-AI^*c_1\sqrt{\hat\mu}+b\hat\lambda-L\hat\lambda^2-\tfrac12B_3^*c_1\sqrt{\hat\mu}\,\hat\lambda-\tfrac12\big(P^*\hat\mu+Q^*\big)c_1^2\ \ge\ A-AI^*c_1+b\hat\lambda-L\hat\lambda^2-\tfrac12B_3^*c_1\hat\lambda-\tfrac12(P^*+Q^*)c_1^2=\sigma^*(\hat\lambda,c_1;1),
\]
using $\sqrt{\hat\mu}\le1$ and $\hat\mu\le1$ in the negative terms; the corners with $e=0$ give $A>0$ and $A+b\lambda_1-L\lambda_1^2>0$ ($\lambda_1\le1<\lambda_c$). So it suffices that $\sigma^*(0,c_1;1)\ge0$ and $\sigma^*(\lambda_1,c_1;1)\ge0$ for each listed pair, which is how the $c_1$ were determined (bisection on $c_1$, rounded down); explicitly, $\sigma^*(0,c_1;1)=A(1-I^*c_1)-\frac12(P^*+Q^*)c_1^2=3.8\cdot10^{-5}$ at $c_1=0.3014$, and $\sigma^*(\lambda_1,c_1;1)$ equals $1.2\cdot10^{-4}$, $4.1\cdot10^{-5}$, $1.1\cdot10^{-4}$, $2.3\cdot10^{-4}$, $6.3\cdot10^{-4}$ at the five pairs with $\lambda_1>0$, each positive.

\emph{(v).} \eqref{eq:s14:struct}, $\hat\lambda S_{01,k}\ge0$ (\cref{lem:s14:A}(i)) and (i).

\emph{(vi).} At $e=0$ only the orders $(0,1)$ and $(0,2)$ of \eqref{eq:s14:CTdef} are present, so by \cref{lem:s14:A}(iii), \cref{lem:s14:C}(iii) (for an arbitrary $M_1\in(0,m)$, and now in general units) $\lim_ks^\infty_k(M_1)=g_\lambda(M_1)\defeq(m-M_1)\big[\frac{m+M_1}2+\frac{(N+2)\lambda}{2\pi}\big]-\frac{(N+2)\lambda^2}{\pi^2}\ln\frac m{M_1}$ exactly, and $\sigma_1(M_1)=g_\lambda(M_1)$. As a function of $M_1\in(0,m]$: $g_\lambda(m)=0$, $g_\lambda'(M_1)=-M_1-\frac{(N+2)\lambda}{2\pi}+\frac{(N+2)\lambda^2}{\pi^2M_1}$, $g_\lambda''(M_1)=-1-\frac{(N+2)\lambda^2}{\pi^2M_1^2}<0$, so $g_\lambda$ is strictly concave. If $g_\lambda'(m)\ge0$, concavity gives $g_\lambda(M_1)<g_\lambda(m)+g'_\lambda(m)(M_1-m)\le0$ for every $M_1<m$, and \Tthree\ fails; if $g'_\lambda(m)<0$, then $g_\lambda>0$ on a left neighbourhood of $m$, and \Tthree\ holds. Now $g'_\lambda(m)<0$ if and only if $x=\lambda/m$ lies below the positive root of $\frac{N+2}{\pi^2}x^2-\frac{N+2}{2\pi}x-1$, which is $x_*=\frac\pi4\big(1+\sqrt{1+\frac{16}{N+2}}\big)$, equal to $\frac\pi4(1+\sqrt5)=2.5416\ldots$ for $N=2$. At a fixed $M_1$, $g_\lambda(M_1)>0$ if and only if $L\lambda^2-b\lambda-A<0$ with $A,b,L$ at $M_1$, i.e.\ $\lambda<\lambda_c(M_1)m\defeq\frac{b+\sqrt{b^2+4AL}}{2L}m$; at $M_1=0.7m$, $\lambda_c=2.14398\ldots$.

\emph{(vii).} By \cref{thm:a06:CD1}(b) the relative limit $\alpha_\infty(M)-\alpha_\infty(m)$ exists for every $M>0$ and $\sigma_1(M_1)=\lim_ks^\infty_k(M_1)=\frac12(\alpha_\infty(m)-\alpha_\infty(M_1))$. On each region of (iii), (iv), (vi), $\sigma_1(0.7m)>0$ (resp.\ $\sigma_1(M_1)>0$), so $\alpha_\infty(0.7m)<\alpha_\infty(m)$: $m$ is not a minimiser of $\alpha_\infty$ on $(0,m]$, which is \sxivNEl, and \sxivNEl\ implies \sxivNE\ (\cref{lem:s14:CW}(b)(i)). \Cref{prop:s10:Upm} is needed only on $E^<_W\setminus E_W$, which these regions avoid.
\end{proof}


\begin{remark}[the quantifier over the couplings]\label{rem:s14:quantifier}
The theorems of this paper print the sources' quantifier of \cref{lem:s14:CW}(e), in the units of the coarsest spacing, with the cube family $\CJ$ not rescaled and $\beta=5^{-J_{\rm K}}$. A sentence ``$e,\lambda$ sufficiently small, $m,\lmu>0$ fixed'' describes \emph{where} the smallness lives (the running couplings, through Ba{\l}aban's $\eps_0$), but as a hypothesis on the bare couplings it is strictly stronger than what the sources need. (i) All four couplings have mass dimension one, so its dimensionless content is a bound on $e^2/m$ and $\lambda/m$, while the sources' hypothesis on the bare couplings is empty beyond positivity: the point $(e^2,\lambda,m,\lmu)=(100,100,1,1)$ is admitted by every source (with a large $J_{\rm K}$) and excluded by the phrase in every system of units. (ii) In the units in which the stopping scale is $1$, $m$ and $\lmu$ are not fixed but rescaled to $m\beta$ and $\lmu\beta$, of which $m\beta$ is capped (by the mass cap of \cite[Prop.~3.1]{Ba2} and by a further cap of the schedule) and $\lmu\beta$ is not: in those units the theorem reads ``three couplings small, the fourth rescaled''. (iii) Even at a fixed stopping scale $\beta_0$ the admissible set is not a product neighbourhood of $(0,0)$ in $(e,\lambda)$: one cap of the schedule of constants of the rotation part has the form $e\le\lambda^{1/4}c(\beta_0)$, a cusp that fails as $\lambda\to0$ at fixed $e$. So the phrase names no set exactly, and the sources' form survives because $\beta$ can be taken to $0$.
\end{remark}
